% !TeX program = pdflatex
% =============================================================================
%  UOS — Unified Oral Scene. Format Specification v0.2
%  Single-file LaTeX source. Compiles on Overleaf with pdfLaTeX, no extra files.
% =============================================================================
\documentclass[11pt,a4paper]{article}

\usepackage[utf8]{inputenc}
\usepackage[T1]{fontenc}
\usepackage{lmodern}
\usepackage[english]{babel}
\usepackage[margin=2.4cm]{geometry}
\usepackage{amsmath}          % \text in the SH formula
\usepackage{microtype}
\usepackage{booktabs}
\usepackage{longtable}
\usepackage{array}
\usepackage{xcolor}
\usepackage{listings}
\usepackage{enumitem}
\usepackage{fancyhdr}
\usepackage{titlesec}
\usepackage{indentfirst}

% ---------------------------------------------------------------- palette ----
% Defined BEFORE hyperref, whose options below name these colours.
\definecolor{uosblack}{HTML}{000000}
\definecolor{uosgrey}{HTML}{5A6570}
\definecolor{uoslight}{HTML}{F2F5F7}
\definecolor{uosrule}{HTML}{C8D2D9}
\definecolor{uosred}{HTML}{8C2F2F}
\definecolor{uosgreen}{HTML}{2F6B45}
\definecolor{uosamber}{HTML}{8A6212}

\usepackage[colorlinks=true,linkcolor=uosblack,urlcolor=uosblack,citecolor=uosblack]{hyperref}

% ------------------------------------------------------------------ misc ----
% Identificadores como `extensions_required` no se parten, y en una columna estrecha se
% salen al margen. Esto deja a TeX estirar el espaciado en la linea desesperada en vez
% de desbordarla: una linea algo suelta se lee, una que invade el margen no.
\setlength{\emergencystretch}{3em}

\setlist{itemsep=1pt,topsep=3pt,parsep=0pt}
\renewcommand{\arraystretch}{1.25}
\titleformat{\section}{\Large\bfseries\color{uosblack}}{\thesection}{0.7em}{}
\titleformat{\subsection}{\large\bfseries\color{uosblack}}{\thesubsection}{0.7em}{}
\titleformat{\subsubsection}{\normalsize\bfseries}{\thesubsubsection}{0.7em}{}
\pagestyle{fancy}\fancyhf{}
\fancyhead[L]{\footnotesize\color{uosgrey}UOS: Unified Oral Scene v0.2}
\fancyhead[R]{\footnotesize\color{uosgrey}\thepage}
\renewcommand{\headrulewidth}{0.4pt}

% conformance keywords
\newcommand{\MUST}{\textcolor{uosred}{\textbf{MUST}}}
\newcommand{\MUSTNOT}{\textcolor{uosred}{\textbf{MUST NOT}}}
\newcommand{\SHOULD}{\textcolor{uosamber}{\textbf{SHOULD}}}
\newcommand{\SHOULDNOT}{\textcolor{uosamber}{\textbf{SHOULD NOT}}}
\newcommand{\MAY}{\textcolor{uosgreen}{\textbf{MAY}}}
\newcommand{\code}[1]{\texttt{\small #1}}

% Callout: rules above and below rather than a box, so it may break across pages
% and may contain lists and verbatim listings. No tcolorbox/mdframed dependency.
\newenvironment{callout}[1]
 {\par\medskip\noindent{\color{uosblack}\rule{\linewidth}{0.9pt}}\par\nobreak\vspace{3pt}
  \noindent{\bfseries\color{uosblack}#1}\par\nobreak\vspace{2pt}}
 {\par\vspace{2pt}\noindent{\color{uosrule}\rule{\linewidth}{0.6pt}}\par\medskip}

% ------------------------------------------------------------- listings -----
% Only language-level keys here; presentation keys live in the styles below.
\lstdefinelanguage{json}{
  sensitive=false,
  morestring=[b]",
  morecomment=[l]{//},
  literate=
   *{\{}{{{\color{uosblack}\{}}}{1} {\}}{{{\color{uosblack}\}}}}{1}
    {[}{{{\color{uosblack}[}}}{1}   {]}{{{\color{uosblack}]}}}{1},
}
\lstdefinestyle{box}{
  basicstyle=\ttfamily\scriptsize,
  backgroundcolor=\color{uoslight},
  frame=single, rulecolor=\color{uosrule},
  framexleftmargin=4pt, framexrightmargin=4pt,
  xleftmargin=4pt, xrightmargin=4pt,
  breaklines=true, columns=fullflexible, keepspaces=true,
  captionpos=b, aboveskip=8pt, belowskip=8pt,
}
\lstset{style=box}
\lstdefinestyle{py}{style=box,language=Python,
  keywordstyle=\color{uosblack}\bfseries,
  commentstyle=\color{uosgrey}\itshape,
  stringstyle=\color{uosgreen}}
\lstdefinestyle{js}{style=box,language=json,
  showstringspaces=false,
  stringstyle=\color{uosgreen},
  commentstyle=\color{uosgrey}\itshape}
\lstdefinestyle{plain}{style=box,language={}}

% table column helpers
\newcolumntype{L}[1]{>{\raggedright\arraybackslash}p{#1}}
\newcolumntype{C}[1]{>{\centering\arraybackslash}p{#1}}

% =============================================================================
\begin{document}

\begin{titlepage}
\centering
\vspace*{2.5cm}
{\color{uosgrey}\large Format Specification}\\[0.6cm]
{\color{uosblack}\Huge\bfseries UOS: Unified Oral Scene}\\[0.5cm]
{\color{uosblack}\LARGE Version 0.2}\\[1.4cm]
{\large An open container for a complete dental case:\\
geometry, volume, appearance, imaging, clinical findings\\
and the \emph{relations between them}, in one file.}\\[2.2cm]

\begin{tabular}{@{}rl@{}}
\textbf{Status} & Draft; reference implementation complete and shipping\\
\textbf{Container media type} & \code{application/vnd.histora.uos} (proposed, unregistered)\\
\textbf{File extension} & \code{.uos}\\
\textbf{Manifest schema} & \code{schemas/uos-manifest-0.2.schema.json}\\
 & (JSON Schema 2020-12; see \S\ref{sec:artifacts} for its \code{\$id})\\
\textbf{Dependencies} & glTF 2.0 and \code{KHR\_gaussian\_splatting} (Release Candidate)\\
\textbf{References} & DICOM PS3.10; FHIR R4; ISO 3950 (FDI)\\
\textbf{Audience} & Implementers of readers, writers, viewers and agents\\
\end{tabular}

\vspace{1.0cm}

\begin{tabular}{@{}rl@{}}
\textbf{Editor} & Luis Garbayo Fern\'andez\\
\textbf{Copyright} & Luis Garbayo Fern\'andez\\
\textbf{Licence} & Apache License 2.0, covering \textbf{this document and the}\\
 & \textbf{reference implementation alike}. The licence in the repository\\
 & root applies to both; neither is more freely reusable.\\
\textbf{Extension prefix} & \code{histora\_}, the single prefix this specification\\
 & reserves (\S\ref{sec:extensions}).\\
\textbf{Produced under} & Version 0.2 was written during the Becas de Verano ANFAIA\\
 & 2026, on a clinical case contributed by Histora.\\
\textbf{Version} & 0.2 --- last revised \today\\
\end{tabular}

\vspace{0.8cm}
{\color{uosgrey}\small Every figure in this document was measured on a real clinical case
contributed by Histora. This is a versioned, living specification, not the deliverable
of a fixed term. Later versions revise it under the process of
\S\ref{sec:limits}; \S\ref{sec:artifacts} names where the schema for each
version lives.}

\vfill
{\color{uosgrey}\small This document is the contract. An implementation that follows it
can open, render, verify, extend and re-emit a \code{.uos} without access to the
reference implementation and without asking its authors anything.}
\vspace{1.5cm}
\end{titlepage}

\tableofcontents
\newpage

% =============================================================================
\section{Introduction}\label{sec:intro}

\subsection{What a \code{.uos} is}
A dental case today is three unrelated things: one or two STL files from the intraoral
scanner, a directory of several hundred DICOM slices from the CBCT, and a PDF report.
Each lives in its own coordinate system, each opens in a different program, and the
\emph{relations} between them, where this CBCT sits with respect to this scan, which
tooth this finding refers to, live in someone's head or inside proprietary software.
Those relations do not travel with the data.

A \code{.uos} is those three things \textbf{plus the relations between them}, in one
file, addressable over HTTP, verifiable byte by byte, and openable in a browser.

\textbf{The one-sentence definition}. A \code{.uos} is an \textbf{uncompressed ZIP} whose first entry is a \code{manifest.json} that declares every payload it carries by content hash, the named coordinate frames those
payloads live in, and the measured transforms that relate those frames.

\subsection{What a \code{.uos} is not}

\begin{itemize}
  \item \textbf{It is not a pixel format.} UOS never re-encodes source data. DICOM travels
  byte-identical; the STL travels byte-identical. What UOS adds is composition.
  \item \textbf{It is not a replacement for DICOM, glTF or FHIR.} It references them. A
  reader that already has a DICOM parser, a glTF loader and a FHIR client needs almost no
  new code, see \S\ref{sec:cookbook}.
  \item \textbf{It is not a clinical validation.} Everything produced by a model is
  isolated under \code{derived/} and marked \code{investigational} (\S\ref{sec:derived}).
  \item \textbf{It is not yet a standard.} It is a container proposal with a schema
  and a validator, both distributed with the reference implementation. Until a second implementer opens one, it is nothing more.
\end{itemize}

\subsection{How to read this document}\label{sec:howtoread}

This document is long because it is exhaustive, not because all of it is for everybody.
Find your task below and read only that path; each one is self-contained and ends in a
worked recipe.

\begin{itemize}[leftmargin=1.2em]
  \item \textbf{To show a \code{.uos} on screen} --- read \S\ref{sec:container} (the
  container), \S\ref{sec:manifest} (the manifest), \S\ref{sec:scenes} (scenes,
  \emph{in full}) and \S\ref{sec:views} (views), then the recipe \emph{Render it}
  (\S\ref{rec:render}). A minimal viewer is about thirty lines.

  \item \textbf{To open one and check it is valid} --- read \S\ref{sec:container} (the
  container), \S\ref{sec:model} (the data model), \S\ref{sec:manifest} (the manifest) and
  \S\ref{sec:conformance} (the validation algorithm), then the recipe \emph{Open a case
  and verify it} (\S\ref{rec:open}).

  \item \textbf{To add data and re-emit it} --- read \S\ref{sec:model} (the data model),
  \S\ref{sec:manifest} (the manifest), \S\ref{sec:provenance} (versioning: a \code{.uos}
  is never edited in place) and, if the base format has no place for what you carry,
  \S\ref{sec:extensions} (extensions). The recipes are in \S\ref{sec:cookbook}.

  \item \textbf{To act on one as an autonomous agent} --- \S\ref{sec:conformance} and
  Appendix~\ref{app:cheat} are enough to operate: what to verify, what to refuse, what may
  be added and where. Everything else is justification for decisions you do not need to
  re-derive.

  \item \textbf{To judge the design rather than implement it} --- read \S\ref{sec:intro}
  for the problem this format exists to solve, the \emph{Rationale} paragraphs and the
  highlighted boxes for why each decision went the way it did, and \S\ref{sec:limits} for
  what v0.2 deliberately does not do.
\end{itemize}

\subsection{Conformance language}

The key words \MUST, \MUSTNOT, \SHOULD, \SHOULDNOT{} and \MAY{} are to be interpreted as
described in RFC 2119. Sections and paragraphs marked \emph{Rationale} are informative;
everything else is normative.

Two failure classes are used throughout and they are not interchangeable:

\begin{longtable}{L{2.6cm} L{12.8cm}}
\toprule
\textbf{Class} & \textbf{Meaning for a reader}\\
\midrule
\endhead
\textbf{Error} & The container is invalid. A reader \MUST{} refuse to present it as a
valid case. Errors are always about internal contradiction --- a declared hash that does
not match, an asset that is referenced and absent.\\
\textbf{Warning} & The container is valid but carries a weaker guarantee than it could,
and the reader \MUST{} surface that to the user rather than swallow it. An unverified
registration and an externalised original are the two canonical examples.\\
\bottomrule
\end{longtable}

\textbf{The design rule behind every warning in this document}. 
The format's central commitment is that \textbf{two things that are not the same must
never look the same}. An automatic alignment nobody reviewed and one signed by a clinician
are different facts. A DICOM series carried inside the container and one merely referenced
by hash are different guarantees. Where the format cannot enforce the difference, it
\MUST{} at minimum name it.

\subsection{Normative artifacts}\label{sec:artifacts}

\begin{longtable}{L{4.6cm} L{10.8cm}}
\toprule
\textbf{Artifact} & \textbf{Location and role}\\
\midrule
\endhead
Manifest JSON Schema & JSON Schema 2020-12, at
\code{schemas/uos-manifest-0.2.schema.json} in the reference repository and retrievable at
its \code{\$id},
\url{https://raw.githubusercontent.com/ANFAIA/Agentic-Smart-Health/uos-spec-v0.2/schemas/uos-manifest-0.2.schema.json}
--- pinned to a tag, so the identifier always names the same document. It is
\emph{generated} from the reference implementation's types, never hand-written. \textbf{What it does and does not check.} It is the only artifact a third party needs
to check a manifest's \emph{shape}. It cannot check truth: not that a \code{sha256} matches
its asset, that the frame graph reaches the canonical frame, or that \code{derived/} carries
its sidecar. That is why \S\ref{sec:conformance} requires both the schema and the validation
algorithm, and why passing the schema alone is not conformance.\\
Reference implementation & Python package \code{uos}. \textbf{Its public API --- everything
an integrator types --- is English}: \code{read\_manifest}, \code{validate},
\code{write\_uos}, \code{Manifest}, \code{Registration}, \code{report.valid},
\code{report.errors}. The module \emph{files} are not, and deliberately: they are a map of
this implementation, not an API. Nobody imports them --- the package re-exports everything
--- and renaming eleven files to make a listing monolingual is a breaking change for
something no reader types. They map one-to-one onto the sections here:
\code{manifiesto.py} (\S\ref{sec:manifest}), \code{contenedor.py}
(\S\ref{sec:container}), \code{escena.py} (\S\ref{sec:scenes}), \code{volumen.py}
(\S\ref{sec:volume}), \code{clinico.py} (\S\ref{sec:clinical}), \code{derivados.py}
(\S\ref{sec:derived}), \code{vistas.py} (\S\ref{sec:views}), \code{procedencia.py}
(\S\ref{sec:provenance}), \code{validador.py} (\S\ref{sec:conformance}).\\
The measured container & \code{another\_patient.uos} --- 12 entries, 18 views, 3 declared
extensions, a provenance chain of 3 versions. Every real value quoted in this document is
read from it. \textbf{It is not distributed with this specification}: it holds one
patient's case, pseudonymised but real, and it is held by the project rather than published.
A reader should not go looking for a file that does not accompany the document. A
distributable conformance fixture is listed as pending in \S\ref{sec:limits}.\\
\bottomrule
\end{longtable}

Where this document and the reference implementation disagree, \textbf{this document is
wrong and \MUST{} be corrected}; but a reader encountering the disagreement \SHOULD{}
follow the implementation, because that is what the containers in the wild were written
by.

% =============================================================================
\newpage
\section{The container}\label{sec:container}

\subsection{Physical format}

A \code{.uos} file \MUST{} be a ZIP archive (PKZIP, ISO/IEC 21320-1 profile) in which:

\begin{enumerate}
  \item Every entry \MUST{} use compression method \code{STORE} (0). A reader \MUST{}
  report an error on any deflated entry.
  \item The \textbf{first physical entry} \MUST{} be named \code{manifest.json}.
  \item Entry names \MUST{} be relative, use \code{/} as separator, be pure ASCII, and
  \MUSTNOT{} contain a \code{..} path segment or a leading \code{/}.
  \item An entry name ending in \code{/} denotes a \emph{directory asset} (\S\ref{sec:dirasset}).
\end{enumerate}

\paragraph{Rationale: why STORE.} With STORE and the ZIP central directory at the end,
an HTTP client issuing \emph{range requests} reads the index and pulls a single asset
without downloading the case. On the example container that is the difference between
fetching 17.9\,MB of scene and fetching 183\,MB of case. Deflating the archive would take
that away: a compressed entry has to be inflated from its start, so there is no useful
byte range to ask for, and the container stops being addressable and becomes a package you
open whole before touching anything.

\paragraph{And what STORE costs, measured rather than assumed.} It is not free, and the
usual justification --- ``the payloads are already compressed'' --- is only half true here.
On the example container the entries total 183.2\,MB stored against 37.3\,MB deflated: the
\code{.glb} is near its limit at 81\,\%, but the \code{.ply} layers, which are raw
\code{float32}, fall to 13--15\,\%. The container is therefore about five times the size a
compressed archive would be, and a writer should know that before quoting a transfer time. The trade is still taken, because \textbf{the right place to compress is inside the asset,
not in the container}. A Gaussian layer stored as SPZ, or a mesh carrying Draco, shrinks by
the same order \emph{and} stays addressable by range; deflating the ZIP buys the same bytes
and gives up the one property the container exists to provide. Compression belongs to the
payload's own format, where a reader that understands the payload can undo it, rather than
to the envelope, where undoing it is a precondition for reading anything at all.

\paragraph{The precedent, and how far it goes.} \code{.usdz} made the same trade: it is a
ZIP whose entries are stored uncompressed, so a runtime can read a texture in place instead
of unpacking the archive first. It goes one step further and \emph{aligns} each entry,
which is what makes the bytes mappable into memory without a copy. UOS does not align: an
HTTP range request has no such requirement, and the property cannot be granted
retroactively --- alignment moves offsets, so containers written today would not acquire it
when a later version asked for it. A future revision that wants zero-copy local reads has
to add the rule and accept that it splits the format.

\paragraph{Rationale: why the manifest goes first.} It is the format's positive
identification. A reader opens the first bytes, sees the entry and its \code{uos\_version}
field, and knows what it is holding without trusting the file extension or a magic number
we would have had to invent.

\subsection{Media type and identification}

\begin{itemize}
  \item \textbf{Media type:} \code{application/vnd.histora.uos} --- \textbf{proposed}. The \code{vnd.} tree is to be registered with IANA when this specification is published; a vendor-tree registration names an organisation and does not require owning a matching domain. Until it is registered a reader \MUSTNOT{} rely on it and \MUST{} identify by rule 2 above.
  
  \item \textbf{Schema \code{\$id}:} it resolves, and it deliberately does not name a domain nobody holds. An earlier draft pointed at \code{histora.dev}, which is unregistered: a \code{\$id} need only be a URI, but writing \code{https://} tells a reader to fetch it, and fetching failed. It now points into the public repository at a \textbf{tag}, not a branch --- an identifier that returns different content depending on the day does not identify. If the format later acquires a domain of its own, moving the \code{\$id} is a major-version change and nothing else.
  
  \item \textbf{Extension:} \code{.uos}
  
  \item \textbf{Magic:} None of its own. \code{PK\textbackslash x03\textbackslash x04} followed by the local file header of \code{manifest.json}.
  
  \item \textbf{Fragment identifier:} \code{caso.uos\#view=<view.id>} selects a stored view (\S\ref{sec:views}). Readers that do not implement views \MUST{} ignore the fragment.
\end{itemize}

\subsection{Reserved paths}\label{sec:paths}

The following paths are reserved by this specification. A writer \MUSTNOT{} place unrelated
content at them; a reader \MAY{} rely on their meaning.

\begin{longtable}{L{5.4cm} L{3.2cm} L{6.2cm}}
\toprule
\textbf{Path} & \textbf{Required} & \textbf{Contents}\\
\midrule
\endhead
\code{manifest.json} & \MUST, first entry & \S\ref{sec:manifest}\\
\code{scene/} & --- & Mesh scene, Gaussian layers, their descriptors (\S\ref{sec:scenes})\\
\code{scene/scene.glb} & by convention & The glTF binary scene\\
\code{scene/*.gs.json} & --- & Gaussian layer descriptors, extension
\code{ash\_gs\_measured}\\
\code{volume/<id>.volume.json} & --- & The volume sidecar --- dimensions, spacing,
orientation, encoding (\S\ref{sec:volume}). The DICOM series it describes is referenced,
not carried, so this file is normally the only thing under \code{volume/}\\
\code{clinical/observations.json} & --- & Per-tooth clinical layer, extension
\code{ash\_clinical} (\S\ref{sec:clinical})\\
\code{derived/} & --- & \textbf{Everything} produced by inference (\S\ref{sec:derived})\\
\code{views.json} & --- & Stored presentation states (\S\ref{sec:views})\\
\code{provenance/chain.json} & --- & Version chain (\S\ref{sec:provenance})\\
\code{provenance/signatures/} & --- & Reserved for Ed25519 signatures; \textbf{not
specified in v0.2}\\
\bottomrule
\end{longtable}

\textbf{A referenced original occupies no path at all}. 
Acquired originals are named by content address and never by location
(\S\ref{sec:external}), so no reserved path corresponds to one. There is no \code{images/}
holding the photographs and no \code{volume/ct\_001/} holding the slices: those assets are
in the manifest, addressed as \code{sha256:<hex>}, and nowhere in the ZIP. That is the property worth protecting. A path would have to be either the local one --- and
the local path of a clinical case carries the patient's own directory name --- or an invented
logical one, which would be a promise about a location inside an archive the file is not in.
A hash is neither: it says which file it is and nothing about where anyone keeps it.
\newline

\begin{lstlisting}[style=plain, caption={Directory layout of the example container}]
manifest.json                 <- first entry, always
scene/scene.glb               <- glTF scene: mesh + appearance splats
scene/field.ply               <- Gaussian field of the CBCT (measured density)
scene/field.gs.json           <- its column descriptor
scene/composite.ply           <- scanner + CBCT in one frame
scene/composite.gs.json
scene/appearance.gs.json      <- descriptor of the splats embedded in scene.glb
views.json                    <- 18 stored camera states
clinical/observations.json    <- per-tooth findings, keyed by FDI
derived/seg_teeth.bin         <- int16 FDI label per vertex  (layer 3)
derived/seg_teeth.meta.json   <- which model produced it     (layer 3)
provenance/chain.json         <- 3 versions of this case
\end{lstlisting}

% =============================================================================
\section{Data model}\label{sec:model}

\subsection{The three planes}\label{sec:planes}

Everything in a \code{.uos} belongs to exactly one of three planes, and the separation is
deliberate:

\begin{longtable}{L{2.4cm} L{3.2cm} L{9.6cm}}
\toprule
\textbf{Plane} & \textbf{Where} & \textbf{What it is, and what may be done to it}\\
\midrule
\endhead
\textbf{Data} & \code{assets[]} & Acquired measurements --- regulatory layer 1 --- and
the representations UOS computes from them --- regulatory layer 2 --- each declaring its
source in \code{derived\_from} and, where the conversion loses something, saying so
(\S\ref{sec:scenes}). Immutable: a writer never edits an asset in place
(\S\ref{sec:provenance}).\\
\textbf{Presentation} & \code{views.json} & Camera states, layer visibility, clipping.
Layer 1, disposable, versionable, cheap to rewrite.\\
\textbf{Inference} & \code{derived/} & Anything a model produced. Regulatory layer 3.
\textbf{Removable}: deleting the directory and its manifest entries leaves a valid,
complete \code{.uos}.\\
\bottomrule
\end{longtable}

\paragraph{The three regulatory layers.}
\code{regulatory.layer} takes one of three values, and until this revision only two of
them were defined --- so everything computed travelled as layer 1 without anybody saying
so.

\begin{longtable}{L{1.4cm} L{13.9cm}}
\toprule
\textbf{Layer} & \textbf{What it is}\\
\midrule
\endhead
\textbf{1} & \textbf{Acquired}, or transcribed from a document a person signed. A DICOM
series, a scanner mesh, a photograph; the pH someone measured and wrote down.\\
\textbf{2} & \textbf{Computed} by a deterministic, reproducible procedure from layer-1
assets, \emph{with no trained model}. Automatic registrations, format conversions,
subsampling, per-tooth colour measured off a photograph. A layer-2 asset \MAY{} live
outside \code{derived/}, but it \MUST{} declare \code{derived\_from}: if the claim is that
it is reproducible, the container has to say from what, or the claim cannot be checked.\\
\textbf{3} & \textbf{Model output.} Lives only under \code{derived/}
(\S\ref{sec:derived}).\\
\bottomrule
\end{longtable}

Layer 2 is not a softer layer 3. The distinction that matters to a regulator is whether a
trained model is in the path, and the distinction that matters to a reader is whether the
value can be recomputed. Layer 2 answers yes to the second and no to the first, and it is
exactly the class that was invisible before: an ICP transform, an STL converted to glTF, a
subsampled density field, all of them shipping as though they had been measured.

\paragraph{Rationale.} The third plane is what makes the format shippable. Distributing
the same case in a jurisdiction where the AI module is not cleared requires deleting a
directory, not re-exporting anything. That property only holds if it is enforced --- which
is why \S\ref{sec:conformance} makes ``layer 3 outside \code{derived/}'' an error in both
directions.

\subsection{Frames and registrations}\label{sec:frames}

A \textbf{frame} is a named coordinate system. Every asset declares exactly one. One frame
is the \code{canonical\_frame}: the geometric hub of the case, conventionally the intraoral
scanner's frame. A \textbf{registration} is a first-class object carrying the $4\times4$
homogeneous transform that maps points from \code{source\_frame} into \code{target\_frame},
together with \emph{how it was obtained and how well it fits}.

\paragraph{Normative: frame connectivity.} Treat \code{registrations[]} as undirected edges
over the set of frames. Every frame named by any asset \MUST{} be reachable from
\code{canonical\_frame}. A reader \MUST{} report an error otherwise --- an asset in a frame
that cannot be related to the hub would be placed somewhere arbitrary, and nothing
downstream could detect it.

\paragraph{Conventions within a container.}

\begin{itemize}
  \item Units are \textbf{millimetres}, and this one is checked:
  \code{canonical\_frame.units} \MUST{} be \code{"mm"}.
  \item Handedness is \textbf{right}. It is the default every frame carries and no emitter
  varies it, but \textbf{no validator enforces it} --- a frame declaring otherwise is
  accepted and would mirror everything placed through it. Treat it as a convention a reader
  \SHOULD{} check for itself until a later version makes it normative.
  \item \code{transform\_4x4\_row\_major} is exactly 16 numbers in \textbf{row-major}
  order. glTF stores matrices column-major; a writer emitting this transform into a glTF
  node \MUST{} transpose (\S\ref{sec:gsnodes}). Getting this wrong does not crash --- it
  places the cloud rotated and mirrored, which is worse.
\end{itemize}

\subsection{Visits}

\code{visits[]} names the temporal axis. Every asset and every view declares the visit it
belongs to. Longitudinal comparison, the same case with a baseline and a follow-up in
the same frame, is a matter of two visits and one shared registration, not a second
study.

\subsection{Identity: how anything in a \code{.uos} is named and verified}\label{sec:identity}

\subsubsection{File assets}
\code{sha256} is the SHA-256 of the exact bytes stored in the ZIP entry, lowercase hex,
64 characters. \code{bytes} is their length. A reader \MUST{} verify both.

\subsubsection{Directory assets}\label{sec:dirasset}
A \textbf{directory asset} is one whole multi-file series --- in practice a DICOM study.
What identifies it is that it carries a \code{parts[]} array, one entry per file, with a
relative \code{name}, a \code{sha256} and \code{bytes}. Its \code{uri} follows the same
rule as any other asset: a path ending in \code{/} if the series travels, a content address
if it is referenced (\S\ref{sec:external}), which for an acquired series is always the case.

\paragraph{\code{parts[]} travels even when the series does not.} That is the point of
splitting it out: the container cannot verify a series it does not hold, but whoever holds
it can --- file by file, against these hashes. Referencing gives up custody, not
traceability.

\vspace{6px}
The asset's own \code{sha256} is a \textbf{digest over the part names and hashes}, defined
here because the format needs writer and validator to compute it identically:

\begin{lstlisting}[style=py]
def digest_of_parts(parts):
    h = hashlib.sha256()
    for p in sorted(parts, key=lambda x: x.name):     # sort by name, byte order
        h.update(f"{p.name}\0{p.sha256}\n".encode())  # NAME, NUL, HEX HASH, LF
    return h.hexdigest()
\end{lstlisting}

\paragraph{Rationale.} It runs over names and hashes rather than concatenated bytes for
two reasons: renaming a slice changes the digest --- correct, because the ordering of a
DICOM series is clinical data --- and verifying it costs nothing beyond the manifest,
without re-reading the series itself. Per-file hashes are the difference between reporting
``this series does not match'' and reporting ``slice 214 is corrupt''.

\subsubsection{Referenced originals, and content addressing}\label{sec:external}

\paragraph{Normative: acquired originals are referenced, not carried.} No acquired original
\MUST{} travel inside a \code{.uos} --- not the DICOM series, not the scanner STL, not the
clinical photographs, not the reports. They are declared with \code{external: true} and
named by content address. What a container carries is what UOS composes: the scene, the
Gaussian layers, the views, the clinical layer, and anything under \code{derived/}.
\textbf{This is the format, not a variant of it.} There is no ``full'' packaging to opt
into. A container found to carry an acquired original is a writer defect, not a richer
case: it is duplicated source data whose custody the format never claimed. The photographs
are worth naming explicitly because they are the ones that slipped --- eighteen megabytes
of patient images once travelled inside a container that declared it carried no originals,
because they were tied to a different flag from the STL and the DICOM.

\vspace{6px}
An asset that is declared but not carried sets \code{external} to \code{true}, and its
\code{uri} \MUST{} be a \textbf{content address}, not a path:

\begin{lstlisting}[style=js]
{ "id": "asset.ios", "external": true,
  "uri":    "sha256:abca87e971dab2d6a49d60edfe7f73f95ce62410603df86e72373a8c874197cb",
  "sha256":        "abca87e971dab2d6a49d60edfe7f73f95ce62410603df86e72373a8c874197cb",
  "bytes": 11004334, "media_type": "model/stl", "frame": "frame.ios_master" }
\end{lstlisting}

Normative rules, checked in both directions:
\begin{itemize}
  \item \code{external: true} $\Rightarrow$ \code{uri} matches \code{\^{}sha256:[0-9a-f]\{64\}\$}.
  \item \code{uri} is a content address $\Rightarrow$ \code{external} is \code{true}
  (otherwise the asset would be unlocatable inside the ZIP).
  \item The hash inside the \code{uri} \MUST{} equal the \code{sha256} field.
  \item A reader \MUST{} make plain, once per container, which originals are referenced
  rather than carried: the container declares what each file is and cannot prove it holds
  it. This is a statement about custody, not a defect to be flagged per asset --- since
  every acquired original is referenced, a per-asset warning would fire on every container
  and distinguish nothing.
\end{itemize}

\paragraph{What is lost, and what replaces it.} A referenced original is not in the
container, so ``the DICOM that comes out is byte-identical to the one that went in'' is a
claim this format does not make about a file it does not hold. What it does hold is enough
to make two weaker, useful claims: \textbf{identity} --- whoever has the file can prove it
is the one this case was built from, slice by slice (\S\ref{sec:dirasset}) --- and
\textbf{reconstruction}: the \code{ash\_reversible} extension regenerates an arch mesh from
the scene itself (\S\ref{sec:extensions}). Reversibility here means producing the mesh
again from what travels, not handing back the file that was ingested. And what comes back is not the file that went in: the regenerated mesh carries per-vertex
colour sampled from the patient's own photographs and an FDI code per vertex, neither of
which the ingested STL ever had. It is exported as \code{.ply}, \code{.stl} or \code{.3mf}
--- the colour and the codes surviving in the two formats able to carry them --- so the
practical answer to ``why not just keep the STL'' is that the container hands back an
improved one.

\paragraph{Rationale: why not a path.} A file that does not travel has no location
inside the container, so a path would be a promise about a ZIP it is not in. What remains
true of it is \emph{which file it is}, and that is exactly what a content address states.
It also has a property a path does not: the local path of a clinical case contains the
patient's directory name; a hash contains nothing. Referencing takes files out of the
container, never identities.

% =============================================================================
\section{manifest.json}\label{sec:manifest}

The manifest is the contract. What is not declared here cannot be in a \code{.uos}, the
point being that a reader can \emph{refuse} rather than guess. The schema sets
\code{additionalProperties: false} on every object: an unknown field is an error, and new
information is added through the extension mechanism (\S\ref{sec:extensions}), not by
sprinkling keys.

\subsection{Top level}

\begin{longtable}{L{3.5cm} L{3.0cm} C{1.4cm} L{6.4cm}}
\toprule
\textbf{Field} & \textbf{Type} & \textbf{Req.} & \textbf{Description}\\
\midrule
\endhead
\code{uos\_version} & string & default \code{"0.2"} & Format version. A reader \MUST{}
check it before anything else.\\
\code{case\_id} & string & \MUST & Stable case identifier. Conventionally
\code{urn:uuid:<v5>}; \MUSTNOT{} encode patient identity.\\
\code{created} & date-time & --- & RFC 3339, UTC.\\
\code{generator} & map str$\to$str & \MUST & Who wrote this container, e.g.
\code{\{"name": "...", "version": "0.4.0"\}}.\\
\code{phi\_state} & enum & \MUST & \S\ref{sec:phi}.\\
\code{subject} & \code{Subject} & \MUST & \S\ref{sec:phi}.\\
\code{canonical\_frame} & \code{Frame} & \MUST & The geometric hub. \code{units} \MUST{}
be \code{"mm"}.\\
\code{frames} & \code{Frame[]} & --- & Every \emph{other} frame. The canonical one is not
repeated here.\\
\code{visits} & \code{Visit[]} & --- & Temporal axis.\\
\code{assets} & \code{Asset[]} & --- & The data plane.\\
\code{registrations} & \code{Registration[]} & --- & The relations. \S\ref{sec:frames}.\\
\code{fhir\_map} & map str$\to$\code{FhirResource} & --- & Keyed by asset \code{id}, plus
the reserved key \code{"case"} for the container itself. \S\ref{sec:fhir}.\\
\code{extensions} & map str$\to$\code{Extension} & --- & \S\ref{sec:extensions}.\\
\code{extensions\_used} & string[] & --- & \S\ref{sec:extensions}.\\
\code{extensions\_required} & string[] & --- & \S\ref{sec:extensions}.\\
\code{provenance} & \code{Provenance} & --- & \S\ref{sec:provenance}.\\
\bottomrule
\end{longtable}

\subsection{Asset}

\begin{longtable}{L{3.2cm} L{2.9cm} C{1.4cm} L{6.8cm}}
\toprule
\textbf{Field} & \textbf{Type} & \textbf{Req.} & \textbf{Description}\\
\midrule
\endhead
\code{id} & string & \MUST & Unique within the manifest. Convention \code{asset.<name>};
this is the key used by \code{fhir\_map}, \code{derived\_from} and sidecars.\\
\code{kind} & enum & \MUST & \code{volume}, \code{mesh\_gs\_scene}, \code{image2d},
\code{signal}, \code{derived\_seg}, \code{document}.\\
\code{visit} & string & \MUST & A \code{visits[].id}.\\
\code{uri} & string & \MUST & Internal path, or \code{/}-terminated directory, or a
content address --- which is what every acquired original carries (\S\ref{sec:external}).\\
\code{media\_type} & string & \MUST & IANA type of the payload, e.g.
\code{model/gltf-binary}, \code{application/dicom}, \code{model/stl},
\code{image/jpeg}. Declared for referenced assets too: it says what the file \emph{is},
which stays true whether or not the container holds it.\\
\code{sha256} & hex(64) & \MUST & \S\ref{sec:identity}.\\
\code{bytes} & int $\geq 0$ & \MUST & For a directory asset, the sum over \code{parts}.\\
\code{frame} & string & \MUST & Must be reachable from the canonical frame.\\
\code{acquisition} & object & --- & \code{\{time: date-time|null, device: map str->str\}}.\\
\code{load\_priority} & int & default 30 & Progressive-loading hint; lower loads first.
See table below.\\
\code{regulatory} & \code{Regulatory} & --- & \code{\{layer: 1..3, clearances:
Clearance[]\}}, where a \code{Clearance} is \code{\{jurisdiction, regime, status,
reference|null\}} and \code{status} is one of \code{not\_a\_device},
\code{investigational}, \code{submitted}, \code{cleared}, \code{withdrawn}. An empty
\code{clearances} means \textbf{not declared}, by written definition; a validator \MUST{}
warn for each layer-3 asset that leaves it empty. \S\ref{sec:planes}, \S\ref{sec:derived}.\\
\code{projection} & \code{Projection}\,|\,null & --- & 2D imaging only. \S\ref{sec:image2d}.\\
\code{parts} & \code{Part[]} & --- & Directory assets only, carried or referenced.
\S\ref{sec:dirasset}.\\
\code{sidecar\_uri} & string\,|\,null & --- & A description of the asset that spares the
reader from parsing it. If present it \MUST{} exist in the container --- including when the
asset it describes does not (\S\ref{sec:volume}).\\
\code{derived\_from} & string[] & --- & Asset ids this one was produced from. Purely
additive: a reader that ignores it still opens the case.\\
\code{external} & bool & default \code{false} & \S\ref{sec:external}. The default is
\code{false} because it is a property of the envelope, not a recommendation: every acquired
original sets it to \code{true}.\\
\bottomrule
\end{longtable}

\paragraph{Default load priorities.} \code{mesh\_gs\_scene} 10, \code{image2d} 20,
\code{volume} 30, \code{signal} 30, \code{derived\_seg} 40, \code{document} 50. The mesh
scene comes first because it is what lets a viewer show something useful having fetched
only the manifest and one asset. A writer \MAY{} override a single asset's value --- the
example container gives its Gaussian field 25, ahead of the volume it was fitted to.

\paragraph{On \code{derived\_from}, and why the field exists.} \code{scene/scene.glb}
is the scanner mesh re-indexed into glTF: the same 220{,}085 triangles as the STL. Without
this field the manifest declared it as a free-standing \code{mesh\_gs\_scene} with no link
to the \code{asset.ios} it converts, so a recipient had no way to know it was a conversion
rather than an independent source. The information existed --- inside the GLB, as
\code{extras.uos\_source\_asset} --- but recovering it meant parsing eighteen megabytes of
glTF to learn something the manifest states in one line.

\subsection{Registration}

\begin{longtable}{L{4.0cm} L{3.4cm} C{1.3cm} L{5.6cm}}
\toprule
\textbf{Field} & \textbf{Type} & \textbf{Req.} & \textbf{Description}\\
\midrule
\endhead
\code{id} & string & \MUST & e.g. \code{reg.ct\_to\_ios}.\\
\code{source\_frame} & string & \MUST & Points are mapped \emph{from} here.\\
\code{target\_frame} & string & \MUST & \ldots{}\emph{to} here.\\
\code{transform\_4x4\_row\_major} & number[16] & \MUST & Row-major, mm, right-handed.\\
\code{method} & string & \MUST & How it was computed. Free-form; observed values are
\code{icp\_surface}, \code{auto\_dl}, \code{manual}, \code{fiducial}. It describes the
algorithm, \textbf{not} whether a human stood behind it.\\
\code{rms\_error\_mm} & number\,|\,null & --- & Residual of the fit. \code{null} means
\emph{not reported}, never \emph{zero}.\\
\code{computed} & date-time\,|\,null & --- & When.\\
\code{operator} & string\,|\,null & --- & Who or what produced it. A machine \MUST{} write
\code{auto:<agent>@<version>}; this prefix is what marks the registration as unattended.\\
\code{verified\_by} & string\,|\,null & --- & The human who reviewed it. \code{null} means
\textbf{nobody has}.\\
\code{regulatory} & \code{Regulatory} & --- & As for assets.\\
\bottomrule
\end{longtable}

\paragraph{Normative: provisional registrations.} A registration is \textbf{provisional}
when it was produced without a human, \code{operator} beginning \code{auto:}, and
\code{verified\_by} is \code{null}. A validator \MUST{} warn, and a viewer \MUST{} present
it as unverified. An automatic alignment nobody has looked at and one signed by a clinician
are not the same claim, and the format refuses to let them look alike.

\paragraph{The rule keys on the operator, and it must.} The reference implementation
originally keyed it on \code{method == "auto\_dl"}, and the consequence is instructive: the
only registration it emits declares \code{method: "icp\_surface"}, so the case in the
example container, automatic, unverified, 0.666\,mm of residual, \emph{did not trip
the rule it exists for}. A safeguard written against one algorithm's name stops working the
moment someone uses a different algorithm, which is every time. Who computed it is the
question; how, is a separate field.

\begin{lstlisting}[style=js]
{ "id": "reg.ct_to_ios", "source_frame": "frame.ct_001", "target_frame": "frame.ios_master",
  "transform_4x4_row_major": [0.99786, 0.02798, 0.05902, -89.3628,
                             -0.05533,-0.11815, 0.99145, -22.6927,
                              0.03471,-0.99260,-0.11635,  38.7773,
                              0.0,     0.0,     0.0,       1.0],
  "method": "icp_surface", "rms_error_mm": 0.666,
  "operator": "auto:geometric-fusion-agent@0.2.0", "verified_by": null }
\end{lstlisting}

\subsection{PHI and subject}\label{sec:phi}

\code{phi\_state} is an \textbf{explicit state}, one of \code{identified},
\code{pseudonymized}, \code{anonymized}, \code{quarantined}. Without the field, ``I do not
know whether this carries identifiers'' and ``it does not'' would be indistinguishable,
which is the worst way to handle PHI. Viewers \MUST{} honour it in banners and export
policy.

\code{subject} carries \code{pseudonym}, an optional \code{fhir\_patient} reference, and an
optional \code{consent} (\S\ref{sec:consent}). A patient name \MUSTNOT{} appear anywhere in
a \code{.uos}.

\code{pseudonym} \MUST{} be either a random identifier whose correspondence table lives
outside the container, or an \textbf{HMAC} under a per-issuer key. A plain hash of the
clinical identifier, truncated or not, \MUSTNOT{} be used: the space of identifiers a
clinic issues is small, so an unkeyed hash inverts by dictionary and the pseudonym
pseudonymises nothing.

\paragraph{What \code{phi\_state} cannot say on its own.}
It is a claim about DICOM tags, and \textbf{the container identifies a person without any
tags}. \code{scene/field.ply} is the CBCT density including soft tissue, and a facial
surface is reconstructible from it --- under HIPAA Safe Harbor an ``image comparable'' to a
full-face photograph, under the GDPR a biometric datum in the special category. The
dentition identifies on its own; forensic odontology is the field that does it. \textbf{No
value of \code{phi\_state} makes a \code{.uos} non-personal data}, and a writer \MUSTNOT{}
present one as if it did.

So the field stays, and \code{pseudonymized} and \code{anonymized} \MUST{} additionally
carry a \code{deidentification} block naming what was actually done, in the vocabulary
DICOM PS3.15 Annex E already defines:

\begin{lstlisting}[style=js]
"deidentification": {
  "profile": "DICOM PS3.15 E.1 Basic Application Level Confidentiality Profile",
  "options": ["RetainLongitudinalTemporalInformationModifiedDates",
              "CleanDescriptors", "CleanRecognizableVisualFeatures"],
  "applied_to": ["asset.ct_001", "asset.img_001"],
  "tool": {"name": "...", "version": "...", "sha256": "..."},
  "date_shift_days": null,
  "note": "..." }
\end{lstlisting}

\begin{longtable}{L{3.6cm} L{11.8cm}}
\toprule
\textbf{Rule} & \textbf{Statement}\\
\midrule
\endhead
Say what was done & \code{phi\_state} other than \code{identified} \MUST{} carry
\code{deidentification}. Declaring the state without the measures that produced it is a
claim nobody can check. A validator warns for each \code{volume} or \code{image2d} asset
absent from \code{applied\_to}: not being in the list means nothing was applied to it.\\
The face & If \code{phi\_state} is not \code{identified} and the container carries a
Gaussian layer with \code{measured: true} derived from a volume, \code{options} \MUST{}
include \code{CleanRecognizableVisualFeatures}, and the layer that travels \MUST{} be the
one generated \emph{after} it was applied. Otherwise nothing stops a writer generating the
field from the original CBCT and defacing only the referenced series, which is the copy
nobody looks at.\\
The date shift & \code{date\_shift\_days} \MUST{} be \code{null} unless \code{phi\_state}
is \code{identified}. Shifting every date in a case by the same amount preserves
longitudinality and is the recommended option; publishing \emph{how far} undoes exactly the
measure it claims to apply.\\
The device & \code{acquisition.device} has fixed keys --- \code{manufacturer},
\code{model}, \code{software\_version} --- and \textbf{no serial number}. It was a free
\code{str}$\to$\code{str} map, and a free map in a clinical container ends up holding the
serial, which Safe Harbor lists as a direct identifier alongside the name. A flow that
needs it declares \code{phi\_state: identified}.\\
Dates in general & Every \code{date-time} in the manifest is subject to the policy declared
in \code{deidentification}. Safe Harbor keeps only the year of a date relative to an
individual.\\
\bottomrule
\end{longtable}

\paragraph{The obligation this puts on a writer.} \code{phi\_state} is a claim about the
\emph{case}, and a case includes originals the container does not hold
(\S\ref{sec:external}). DICOM headers carry \code{PatientName}, \code{PatientBirthDate},
\code{InstitutionName} and friends; a writer declaring \code{pseudonymized} \MUST{} have
inspected those tags at ingestion rather than assumed someone else's work, and it is the
writer's last chance to, because once the series is referenced the container cannot be
re-inspected. Note that \code{PatientID} is deliberately \emph{not} on that list: a
correctly pseudonymised export populates it with an opaque token, and rejecting it would
reject perfectly anonymous series.

\vspace{6px}
Referencing does not weaken this and does not strengthen it either. It keeps identifiers out
of \emph{this file} --- a content address carries no patient directory, no filename, no tags
--- but the tags still exist in the series someone else holds, and \code{phi\_state} still
speaks for them.

\subsection{What already exists elsewhere, and where UOS maps to it}\label{sec:mappings}

Four places where this format would otherwise reinvent something a standard already
defines. Naming the destination costs nothing now and saves an incompatible extension
later.

\begin{longtable}{L{4.0cm} L{11.4cm}}
\toprule
\textbf{What UOS has} & \textbf{Where it maps}\\
\midrule
\endhead
\code{registrations[]} & DICOM's \textbf{Spatial Registration Object} (Spatial Registration
Storage, SOP class \code{1.2.840.10008.5.1.4.1.1.66.1}), which references two Frame of
Reference UIDs and carries the matrix --- a PACS can store, version and attach it to the
study. \code{source\_frame}/\code{target\_frame} map to those two UIDs
(\S\ref{sec:frameid}) and \code{transform\_4x4\_row\_major} to the Frame of Reference
Transformation Matrix, which DICOM also stores row-major. In FHIR the registration attaches
to \code{ImagingStudy} with a note.\\
The referenced STL & DICOM defines \textbf{Encapsulated STL Storage} (Supplement 205, SOP
class \code{1.2.840.10008.5.1.4.1.1.104.3}). When the recipient is a PACS, that is the
destination rather than a loose file beside it.\\
\code{provenance/chain.json} & FHIR \code{Provenance}, field for field: which entity,
generated by which agent, when, superseding which. One resource per link, with
\code{entity.role = revision} and \code{target} the manifest's \code{DocumentReference}.\\
\code{projection.fdi\_targets} & FHIR R5's \code{ImagingSelection}, which references a
selection inside an imaging study --- instances, frames, regions. v0.2 declares R4, and
naming R5's destination here is what stops somebody inventing a private FHIR extension for
something R5 already has.\\
\bottomrule
\end{longtable}

\subsection{Clinical vocabulary}\label{sec:vocab}

\code{findings: ["aparato\_ortodoncico"]} interoperates with nothing: not with the FHIR
\code{Observation} it maps to, which expects a coded \code{code}; not with a foreign
practice-management system; not with a second implementer. A closed vocabulary of one
issuer's own terms, in one language, is good internal control and \MUSTNOT{} travel as the
code.

So a finding travels as \code{\{system, code, display\}}. The issuer's own term survives in
\code{display}; \code{code} is SNOMED CT, and for caries the ICDAS 0--6 severity scale.
\code{bodySite} is the FDI/ISO 3950 code.

\paragraph{And \code{code} may be \code{null}, deliberately.}
Assigning a SNOMED identifier is an act of clinical terminology, not of programming: it
needs a licensed browser --- SNOMED CT is licensed per country --- and somebody who answers
for the equivalence. A code guessed at would be indistinguishable, to the connector that
resolves it against a terminology server, from one that is right: the plausible, silent,
already-inside-the-clinical-record failure. \code{null} says \emph{not mapped}; a number
put there by hand says \emph{mapped}, and lies. It is the same rule that leaves
\code{weights\_sha256} null.

\paragraph{FDI is not the only numbering.}
A US reader sees \code{"27"} and reads tooth 27 of the \textbf{Universal Numbering System}
(permanent 1--32, deciduous A--T), which is a different tooth. A container declaring ISO
3950 \MUST{} say so, and an implementer serving both markets needs the mapping written
down rather than assumed.

\subsection{Consent and purpose of use}\label{sec:consent}

A \code{.uos} leaving the clinic for a dental laboratory, for a colleague's second opinion,
and for a model-training pipeline are \textbf{three different legal acts}. Until this
revision the container distinguished none of them, so every recipient had to assume one.
Purpose of use is the first question any data-protection review asks, and it is what
decides which layers may travel at all.

\begin{lstlisting}[style=js]
"subject": { "pseudonym": "59be9708...",
             "consent": { "fhir_consent": "Consent/c-1234",
                          "scope": ["treatment", "lab_manufacturing"],
                          "obtained": "2026-08-01" } },
"purpose_of_use": "lab_manufacturing"
\end{lstlisting}

\code{scope} and \code{purpose\_of\_use} draw from one closed vocabulary:
\code{treatment}, \code{lab\_manufacturing}, \code{second\_opinion}, \code{research},
\code{model\_training}. \code{fhir\_consent} follows \S\ref{sec:fhir}'s rule --- a known
type, a \code{null} reference until the resource exists on a server; FHIR R4's
\code{Consent} is its natural destination.

\code{purpose\_of\_use} declares what \emph{this} container was emitted for and \MUST{} be
contained in \code{subject.consent.scope}: a container cannot be issued for a use the
patient did not consent to. A validator reports an error when it is not, and a warning when
\code{purpose\_of\_use} is absent. A writer \MUSTNOT{} default it to \code{treatment};
assuming a legal act is the failure this field exists to prevent.

\subsection{Occlusion, sites and registration fitness}\label{sec:dentistry}

\paragraph{Occlusion.} Mandible-to-maxilla --- the bite registration --- is the
clinically most important registration of a dental case, and the format did not name it.
The reference case is maxilla-only and that is why it never appeared; that is a reason to
reserve it, not to omit it. The id \code{reg.mandible\_to\_maxilla} is reserved, with
\code{method: "bite\_scan"}. A case with two arches \MUST{} declare either that
registration or \code{occlusion: "not\_recorded"} --- \textbf{silence is never ``there
is none''}. A single-arch case answers \code{single\_arch}, but it answers.

\paragraph{Sites that are not teeth.} \code{uos\_fdi} labels teeth only, so there was no
way to mark an implant site, an edentulous span or a prosthetic abutment --- which is what
half of a rehabilitation is about. \code{uos\_site} is reserved with
\code{implant\_site}, \code{edentulous}, \code{pontic} and \code{abutment}, alongside ISO
3950's area codes for quadrants and arches.

\paragraph{Fitness for use.} \code{rms\_error\_mm} is a global average and \textbf{does not
decide a clinical use}: for guided implant surgery what matters is the maximum local error
in the region of interest, and 0.666\,mm RMS is acceptable for viewing and not for
planning. \code{Registration} therefore carries optional \code{max\_error\_mm},
\code{target\_registration\_error\_mm} --- measured at points \emph{not} used to compute
the registration, unlike the RMS, which is measured on the points that were and is
therefore systematically optimistic --- with the region it was measured in, and
\code{fit\_for}. An empty \code{fit\_for} means \textbf{not declared}, and a reader
\MUSTNOT{} assume fitness for a use that is not listed.

\paragraph{Deciduous dentition.} ISO 3950 covers 51--85 and a validator \MUST{} accept
them.

\subsection{Frame identity and anatomical convention}\label{sec:frameid}

A frame id such as \code{frame.ct\_001} is a string the writer invented. DICOM already
identifies a coordinate system globally and uniquely --- the \textbf{Frame of Reference
UID}, tag \code{(0020,0052)}, which every CBCT series carries --- and a reader that
receives the series through another channel has no way to know that \code{frame.ct\_001}
is \emph{that} series except by trust.

\begin{longtable}{L{3.6cm} L{11.8cm}}
\toprule
\textbf{Rule} & \textbf{Statement}\\
\midrule
\endhead
DICOM identity & \code{Frame} carries \code{dicom\_frame\_of\_reference\_uid}. It \MUST{}
be present on every frame declared by an asset of \code{kind: volume}, the volume sidecar
repeats it, and a validator checks they agree. On writing it \MUST{} be read from the
ingested series and never invented --- the same rule \S\ref{sec:volume} already applies to
\code{orientation}.\\
Anatomical convention & \code{Frame.anatomical} is \code{LPS}, \code{RAS} or
\code{device}. \textbf{``Right-handed'' fixes chirality, not orientation}: DICOM is LPS, an
intraoral scanner uses an arbitrary device system, glTF is Y-up with no anatomical meaning,
and all three can be right-handed while agreeing on nothing useful. Any clinical
measurement over the model --- an angle, a distance to a structure --- needs to know which
direction is the patient's anterior or superior, and until now that meant looking at the
image. Every frame of a \code{volume} \MUST{} declare \code{LPS}, which is what DICOM
imposes.\\
The canonical frame & It \MAY{} be \code{device}, in which case the manifest \SHOULD{}
contain a registration from it to an \code{LPS} or \code{RAS} frame. A validator warns
when neither holds.\\
Frames are per acquisition & Not per device. Two scans of the same patient at different
visits are \textbf{different frames} even from the same scanner. Sharing a frame asserts
that the two clouds are already in one space, and between two visits that is false: the
patient moved, the bite changed, or both.\\
\bottomrule
\end{longtable}

\subsection{FHIR mapping}\label{sec:fhir}

\code{fhir\_map} tells a practice-management connector what to create for each asset.

\begin{longtable}{L{3.4cm} L{2.8cm} C{1.2cm} L{7.4cm}}
\toprule
\textbf{Field} & \textbf{Type} & \textbf{Req.} & \textbf{Description}\\
\midrule
\endhead
\code{resource\_type} & string & \MUST & An R4 resource \emph{type}: \code{ImagingStudy},
\code{DocumentReference}, \code{Media}, \code{Observation}.\\
\code{resource} & string\,|\,null & --- & A \emph{reference} to a resource that exists on
a server, e.g. \code{ImagingStudy/is-9911}. \code{null} until the case actually lives in a
PMS.\\
\code{note} & string & --- & Why this mapping and not another.\\
\bottomrule
\end{longtable}

\paragraph{Rationale, type versus reference.} They are not the same thing and conflating
them is an active hazard: a connector reading a fabricated \code{resource} would try to
resolve it. So a writer declares what it actually knows, the type, and leaves the
reference null.

\paragraph{The honest gap.} FHIR R4 has \textbf{no resource for dental geometry}, and
\code{Media} is defined for photo, video and audio. Meshes and Gaussian layers therefore
map to \code{DocumentReference}, the generic clinical-binary envelope, and the \code{note}
field says so out loud rather than pretending the mapping is exact.

% =============================================================================
\section{Scenes}\label{sec:scenes}

This is the section a viewer implementer needs in full. ``Scene'' in UOS means three
distinct, separable things, and conflating them is the most common way to misread a
container.

\paragraph{First, what a Gaussian layer is here.} A \code{.uos} carries Gaussian layers of
\textbf{two kinds that must never be confused}, and the format spends an entire extension
keeping them apart:

\begin{itemize}
  \item A \textbf{measured} layer is a physical measurement expressed as Gaussians. The
  CBCT density field is one primitive per occupied voxel carrying the attenuation the
  scanner actually measured --- \code{density} is normalised sigma, Beer--Lambert, in the
  patient's frame and in linear millimetres. It is clinical data that happens to be
  rendered by a splat rasteriser, not a rendering that happens to look like data.
  \item A \textbf{reconstructed} layer is appearance fitted by an optimiser: the trained
  colour of the arch. The colour \textbf{is the patient's own}, sampled from their
  photographs; what \code{measured: false} denies is narrower and specific --- the
  optimiser moved, split and pruned the primitives, so no Gaussian corresponds one-to-one
  with the pixels its colour came from. Provenance of the colour and traceability of each
  primitive are different claims, and only the second is lost.
\end{itemize}

De-facto 3DGS knows only the second kind, and its column names --- \code{density},
\code{scale}, \code{opacity} --- carry that meaning by default. Reading a measured layer
with those assumptions produces an image that looks right and means nothing. That is why
every layer travels with a descriptor stating which kind it is (\S\ref{sec:gsjson}), and
why \code{ash\_gs\_measured} exists at all: it is the part of this format the base draft
has no vocabulary for.

\paragraph{And what \code{mesh\_gs\_scene} means.} It is the asset \code{kind}
(\S\ref{sec:manifest}) for anything that is part of the renderable scene --- the glTF
binary and every Gaussian layer alike. One \code{kind} covers both because a viewer treats
them as one thing to compose; what separates a mesh from a splat layer, and a measured
layer from a reconstructed one, is said elsewhere and more precisely.

\vspace{6px}
With that established, the three things:

\begin{longtable}{L{3.6cm} L{4.4cm} L{7.0cm}}
\toprule
\textbf{Thing} & \textbf{Where} & \textbf{What it is}\\
\midrule
\endhead
\textbf{The scene graph} & \code{scene/scene.glb} & A glTF 2.0 binary asset. The node
hierarchy \emph{is} the spatial composition: node 0 is the canonical frame and every other
layer hangs off it with its registration baked into the node transform.\\
\textbf{Embedded appearance} & a primitive inside that GLB & A
\code{KHR\_gaussian\_splatting} point primitive. Any conformant glTF loader renders it
without knowing anything about UOS.\\
\textbf{External Gaussian layers} & \code{scene/*.ply} + \code{*.gs.json} & Layers too
large or too semantically specific to embed --- notably the \emph{measured} CBCT density
field. Referenced from a glTF node via \code{extras.uos\_gs\_uri}.\\
\bottomrule
\end{longtable}

\subsection{Why a glTF conversion exists at all}

An STL is a triangle soup: no indices, no attributes, not one metadata field. Semantic
picking (\S\ref{sec:picking}) is defined over per-sub-mesh metadata, and an STL cannot
carry any. glTF can.

\paragraph{Normative: the conversion is lossy, so the original stays declared.}
glTF requires \code{float32} positions and the ingested mesh is \code{float64}. At
$\pm$100\,mm that leaves roughly 10\,nm of resolution, irrelevant for looking, not
irrelevant for claiming reversibility. Therefore a container carrying
\code{scene/scene.glb} \MUST{} declare the source asset it was converted from --- by
content address, since acquired originals are referenced and not carried
(\S\ref{sec:external}) --- and the GLB \SHOULD{} name it in \code{extras.uos\_source\_asset}
\emph{and} in the manifest's \code{derived\_from}. The converted scene is
\textbf{presentation}; the original is the record, and the record is identified even where
it is not held.

\vspace{6px}

Measured reversibility on the reference case: the mesh regenerates from the container with
$4.59\times10^{-6}$\,mm of error against a 0.1\,mm budget --- an error of the order of
\code{float32}'s rounding unit, which is what the measurement actually shows.

\paragraph{And that is not why the original stays declared.}
This section used to argue that converting to glTF is lossy because ``the ingested mesh is
\code{float64}''. \textbf{That is false for an STL.} A binary STL stores vertices as
\code{float32} \emph{by definition of the format}; the \code{float64} is the loader's array,
a widening of data that was never more precise. Going back to \code{float32} for glTF loses
nothing, and the residual above is the STL's own \code{float32}. From an \textbf{OBJ} ---
decimal text, typically eight digits --- the loss is real, so the claim was true for half
the inputs and stated for all of them.

The rule that \code{scene.glb} \MUST{} declare its source in \code{derived\_from} stands,
for the right reason: \textbf{provenance and regulatory layer} (\S\ref{sec:planes}), not
resolution. A converted mesh is a layer-2 representation of a layer-1 acquisition, and that
is true no matter how exact the conversion is.

\subsection{The scene graph}\label{sec:graph}

\begin{lstlisting}[style=plain]
scene 0
`-- node 0  "scan"                       mesh 0   <- THE CANONICAL FRAME
    |-- node 1  "apariencia real..."     mesh 1   <- KHR_gaussian_splatting primitive
    |                                              (no matrix: trained in the canonical
    |                                               frame, so transform = identity)
    |-- node 2  "campo gaussiano..."     no mesh  <- external layer, extras.uos_gs_uri
    |                                              matrix = reg.ct_to_ios, COLUMN-major
    `-- node 3  "compuesto CBCT+escaner" no mesh  <- external layer
\end{lstlisting}

Normative rules for the graph:

\begin{enumerate}
  \item \code{scenes[0].nodes} \MUST{} contain exactly the root node, and that root node
  \MUST{} carry the mesh geometry expressed in the canonical frame.
  \item Every Gaussian layer, embedded or external, \MUST{} be a \textbf{child} of the root
  node. Its node transform \MUST{} be the registration that carries it into the canonical
  frame. The identity transform \MAY{} be omitted.
  \item Node \code{matrix} is \textbf{column-major} (glTF rule) while
  \code{transform\_4x4\_row\_major} in the manifest is row-major. A writer \MUST{}
  transpose; a reader converting the other way \MUST{} transpose back.
  \item A node transform on a splat node \MUST{} decompose into translation, rotation and
  \emph{positive} scale (a \code{KHR\_gaussian\_splatting} requirement). Registrations from
  rigid alignment satisfy this; a mirrored or degenerate matrix does not, and its rendering
  is undefined.
  \item The manifest remains the auditable statement of the same relation. A reader
  \MAY{} place layers using only the GLB --- which is the point: an off-the-shelf
  \code{GLTFLoader} composes the case correctly with no UOS knowledge --- but an auditor
  \MUST{} read \code{registrations[]}, which is where the error and the verification status
  live.
\end{enumerate}

\subsubsection{Scene-level extras}
\begin{lstlisting}[style=js]
"extras": { "uos_frame": "frame.ios_master", "uos_units": "mm",
            "uos_source_asset": "asset.ios",
            "uos_note": "presentation: float32 from float64. The original is asset.ios,
                         referenced by content address and not carried here; the mesh is
                         regenerated from this scene" }
\end{lstlisting}

\subsection{Mesh primitives and semantic picking}\label{sec:picking}

The mesh travels \textbf{unpartitioned}: a single primitive, a single index buffer, and
\emph{no} \code{uos\_fdi} anywhere. \code{scene/scene.glb} is a layer-1 asset and it
carries layer-1 content only.

\begin{longtable}{L{4.2cm} L{11.2cm}}
\toprule
\textbf{Rule} & \textbf{Statement}\\
\midrule
\endhead
No FDI inside any layer & An asset whose \code{regulatory.layer} is not 3 \MUSTNOT{}
carry the FDI code in \emph{any} of its three forms: \code{extras.uos\_fdi} on a mesh
primitive, the \code{\_REGION\_ID} attribute on a \code{KHR\_gaussian\_splatting}
primitive, or a \code{region\_id} property on a Gaussian \code{.ply}. All three are the
same code and all three come from a segmentation model.\\
Where the codes go instead & \code{derived/seg\_teeth.bin} carries one code per mesh
\emph{vertex}; \code{derived/seg\_gaussians.<layer>.bin} carries one code per
\emph{Gaussian} of the named layer, one file per layer because each has its own count and
its own order. Both are \code{int16} little-endian with a \code{.meta.json} sidecar
declaring what they index.\\
Vertex order is invariant & The scene \MUST{} preserve the vertex order of the mesh it was
built from. This is what allows \code{derived/seg\_teeth.bin} to be a flat array indexed by
vertex, joined by position and nothing else. A writer that re-orders vertices here breaks
that join silently --- and with the scene unpartitioned, that join is now the \emph{only}
route to per-tooth selection.\\
Where picking comes from & A reader that wants per-tooth selection builds it at load time
from \code{derived/seg\_teeth.bin}, by vertex index. For Gaussians it is the same join that
produced them in the first place: nearest crown vertex. \textbf{Picking without
\code{derived/} is not available}, and a writer \MUSTNOT{} make it available by baking the
codes into the scene.\\
Code type & Where an FDI code does appear --- in \code{derived/}, or in a layer-3 asset ---
it is a \textbf{string} (\code{"27"}, not \code{27}). An FDI code is not a quantity, and
writing it as a number invites arithmetic on it.\\
\bottomrule
\end{longtable}

\paragraph{What this costs, stated plainly.}
A foreign glTF viewer opens the container, draws the arch, and \emph{cannot select a
tooth}. That is a real loss of interoperability and it is accepted deliberately.

Version 0.2 of this document specified the opposite: the mesh was split into one primitive
per tooth, each carrying its code, so that picking worked in any glTF viewer. The partition
comes from a segmentation model --- layer 3 --- and it was baked into a layer-1 scene, so
deleting \code{derived/} no longer erased every trace of inference; it left the mesh split
into fourteen pieces carrying their codes. \S\ref{sec:planes} justifies the third plane
with a concrete case: distributing the case ``in a jurisdiction where the AI module is not
authorised'' by deleting a directory. With the partition baked in, that distribution still
contained model output, and no documented compromise sustains the claim in front of an
auditor. The rule is therefore not a writer's option but the only conforming form.

\paragraph{If a generic glTF reader really needs the selection.}
The conforming alternative is a \textbf{second} GLB under \code{derived/}, declared layer
3, with its model sidecar --- never the same file. A reader that accepts \code{derived/}
gets the partitioned scene; a distribution that strips it loses the file and not just its
resolution.

\paragraph{A \code{measured} layer must declare its occupancy threshold.}
``One primitive per \emph{occupied} voxel'' implies a density threshold, and that threshold
is a decision that \textbf{changes which tissue appears}: raising it erases dentine before
enamel. A descriptor that says \code{measured: true} and stays silent about it is applying
the silence rule backwards --- it asserts the field \emph{is} the measurement when it is
the measurement above a cut somebody chose. The \code{histora-twin/1.0} profile therefore
\MUST{} carry \code{occupancy: \{threshold, unit, note\}} and
\code{subsampling.method} (\code{stride}, \code{mean} or \code{max}). Both are reproducible
from the referenced DICOM; without them the layer is not.

\subsection{Appearance: the \code{KHR\_gaussian\_splatting} primitive}\label{sec:khr}

\paragraph{Upstream status.} \code{KHR\_gaussian\_splatting} is a Khronos \emph{Release
Candidate} at the time of writing, written against the glTF 2.0 spec. UOS v0.2 emits it and
additionally keeps the external-layer fallback of \S\ref{sec:gsnodes}. When the extension
is ratified, UOS v1.0 retires that fallback.

\vspace{6px}
The primitive \MUST{} use \code{mode: 0} (POINTS) --- each Gaussian is a point, and the
ellipse is raised by the rasteriser from scale and rotation. The extension object carries
two required properties:

\begin{lstlisting}[style=js]
"extensions": { "KHR_gaussian_splatting": { "kernel": "ellipse",
                                            "colorSpace": "srgb_rec709_display" } }
\end{lstlisting}

\begin{longtable}{L{3.0cm} L{1.6cm} C{3.2cm} L{7.6cm}}
\toprule
\textbf{Property} & \textbf{Type} & \textbf{Req.} & \textbf{Description}\\
\midrule
\endhead
\code{kernel} & string & \MUST & \code{"ellipse"} is the only defined value.\\
\code{colorSpace} & string & \MUST & \code{"srgb\_rec709\_display"} or
\code{"lin\_rec709\_display"} (ASWF CIF naming). \textbf{No sensible default exists};
getting it wrong shifts the hue of the whole arch. UOS appearance layers trained against
sRGB renders declare the former.\\
\code{projection} & string & default \code{"perspective"} & Splat rendering assumes a
perspective camera.\\
\code{sortingMethod} & string & default \code{"cameraDistance"} & Back-to-front sorting.\\
\bottomrule
\end{longtable}

\subsubsection{Attributes}

\begin{longtable}{L{7.0cm} L{1.8cm} C{1.8cm} L{4.6cm}}
\toprule
\textbf{Semantic} & \textbf{Type} & \textbf{Req.} & \textbf{Meaning / units}\\
\midrule
\endhead
\code{POSITION} & VEC3 f32 & \MUST & Gaussian centre, mm, in this node's frame.
\code{min}/\code{max} are mandatory in glTF for this accessor.\\
\code{KHR\_gaussian\_splatting:ROTATION} & VEC4 & \MUST & Unit quaternion in
\textbf{glTF order $(x,y,z,w)$}.\\
\code{KHR\_gaussian\_splatting:SCALE} & VEC3 & \MUST & \textbf{Linear and non-negative},
mm.\\
\code{KHR\_gaussian\_splatting:OPACITY} & SCALAR & \MUST & \textbf{Linear}, in $[0,1]$.\\
\code{KHR\_gaussian\_splatting:}\allowbreak\code{SH\_DEGREE\_0\_COEF\_0} & VEC3 f32 & \MUST & DC term of the
spherical-harmonic colour.\\
\code{KHR\_gaussian\_splatting:}\allowbreak\code{SH\_DEGREE\_$l$\_COEF\_$k$} & VEC3 f32 & \MAY & Degrees
1--3. If a degree is present, \textbf{all lower degrees \MUST{} be present and complete}
(degree 1 is exactly three VEC3 coefficients). Partial degrees are prohibited.\\
\code{COLOR\_0} & VEC3\allowbreak/VEC4 & \SHOULD & Fallback diffuse colour for renderers without the
extension. See below.\\
\code{NORMAL} & VEC3 f32 & \MAY & Standard glTF semantic; a point primitive may carry it.
UOS uses the nearest mesh vertex's normal.\\
\code{\_AO} & SCALAR f32 & \MAY & \textbf{Application attribute} (UOS). Ambient occlusion
in $[0,1]$, to be multiplied into the colour at draw time.\\
\code{\_REGION\_ID} & SCALAR i16 & \MAY & \textbf{Application attribute} (UOS). FDI code
per Gaussian.\\
\bottomrule
\end{longtable}

\paragraph{Reconstructing colour from SH.} The zeroth-order coefficient is not a colour;
it becomes one through the SH basis constant and the training bias:
\[
  \mathrm{RGB}_{\text{diffuse}} \;=\; \mathrm{SH}_{0,0} \times 0.2820947917738781 \;+\; 0.5
\]
A writer emitting the \code{COLOR\_0} fallback \MUST{} compute it as that expression
clamped to $[0,1]$; for an sRGB \code{colorSpace} it must be converted from sRGB to linear
before storage. \textbf{UOS v0.2 does not emit \code{COLOR\_0}}, which means a glTF viewer
without the extension draws the appearance layer as an uncoloured point cloud. Writers
\SHOULD{} emit it; this is expected to become a \MUST{} in v1.0.

\paragraph{Extent.} Splats assume a $3\sigma$ cut-off (Mahalanobis distance 3). This is the
same convention the \code{ash\_reversible} extension uses when sampling colour per vertex,
and it is cited rather than redefined.

\paragraph{Underscore attributes.} \code{\_AO} and \code{\_REGION\_ID} use the underscore
prefix glTF reserves for what its specification does not define. A conformant viewer
ignores them and draws correctly; the UOS viewer uses them for shading and for tooth
selection. That they are non-standard is precisely why the \code{ash\_gs\_measured}
descriptor still exists.

\paragraph{Normative: unit conversions happen \emph{before} this primitive.}
An INRIA-convention 3DGS \code{.ply} stores opacity as a logit, scales as logarithms and
the quaternion as $(w,x,y,z)$. This extension requires linear opacity in $[0,1]$,
non-negative linear scale, and glTF quaternion order $(x,y,z,w)$. A writer \MUST{} apply
sigmoid, exponential and reordering before emitting, and \SHOULD{} validate:
\begin{itemize}
  \item any negative \code{SCALE} means a logarithmic scale was never exponentiated;
  \item any \code{OPACITY} outside $[0,1]$ means a logit never went through the sigmoid.
\end{itemize}
The reference implementation raises on both, because a container that ships either renders
plausibly and wrongly.

\subsection{External Gaussian layers}\label{sec:gsnodes}

Layers that do not embed --- the measured CBCT density field is 1{,}341{,}990 primitives
and 78\,MB --- travel as separate entries inside the container and are referenced from a
child node:

\begin{lstlisting}[style=js]
{ "name": "campo gaussiano del twin",
  "extras": { "uos_gs_uri": "scene/field.ply",
              "uos_descriptor_uri": "scene/field.gs.json",
              "uos_measured": true },
  "matrix": [ /* 16 numbers, COLUMN-major: the transpose of reg.ct_to_ios */ ] }
\end{lstlisting}

A reader that does not understand \code{uos\_gs\_uri} sees an empty node and renders the
rest of the scene correctly. A reader that does \MUST{} read the descriptor
(\S\ref{sec:gsjson}) before interpreting any column, because the column names collide with
de-facto 3DGS and \emph{mean something else}.

\subsection{The \code{.gs.json} layer descriptor}\label{sec:gsjson}

Declared by extension \code{ash\_gs\_measured}. One descriptor per Gaussian layer,
column by column.

\begin{longtable}{L{3.0cm} L{2.2cm} L{9.8cm}}
\toprule
\textbf{Field} & \textbf{Type} & \textbf{Description}\\
\midrule
\endhead
\code{role} & string & Human-readable role of the layer.\\
\code{measured} & bool & \textbf{The load-bearing field.} \code{true} = the values are a
physical measurement; \code{false} = they are a reconstruction.\\
\code{profile} & string & e.g. \code{histora-twin/1.0}, \code{histora-gs-apariencia/1.0}.\\
\code{frame} & string & The frame the coordinates are in.\\
\code{units} & string & \code{"mm"}.\\
\code{n\_primitives} & int & Gaussian count.\\
\code{columns} & object[] & Per column: \code{name}, \code{unit}, \code{scale}
(\code{lineal}\,|\,\code{log}), \code{measured}, \code{derived\_from}, \code{meaning},
\code{vocabulary}.\\
\code{note} & string & What a consumer would otherwise get wrong.\\
\multicolumn{3}{@{}p{14.8cm}@{}}{\small A profile \MAY{} add keys of its own ---
\code{histora-twin/1.0} carries \code{subsampling} with the voxel step and the counts before
and after. A reader \MUST{} ignore keys it does not know; the fields above are the ones
every descriptor has.}\\
\bottomrule
\end{longtable}

\paragraph{Why this descriptor is not optional in practice.}
The CBCT field's \code{density} column is \textbf{normalised sigma --- Beer--Lambert
attenuation --- and not opacity}, and its scales are in \emph{linear millimetres, not
logarithms}. Those column names are identical to de-facto 3DGS and mean something else. A
consumer that assumes the usual convention will exponentiate a value that is already
linear and render a physically meaningless image that looks fine.

\newpage
Two real descriptors, abridged:

\begin{lstlisting}[style=js]
// scene/field.gs.json -- the CBCT density field
{ "role": "campo gaussiano del twin", "measured": true, "profile": "histora-twin/1.0",
  "frame": "frame.ct_001", "units": "mm", "n_primitives": 1341990,
  "note": "density MEASURED by the CBCT: `density` is normalised sigma, not opacity, and
           the scales are in linear millimetres, NOT logarithms",
  "subsampling": {"step_voxels": [3,3,1], "from": 12072785, "to": 1341990},
  "columns": [ {"name":"x","unit":"mm","scale":"lineal","measured":true,  ...},
               {"name":"scale_0","unit":"mm","scale":"lineal","measured":true, ...},
               {"name":"density","unit":"sigma_normalizada","scale":"lineal",
                "measured":true, ...},
               {"name":"region_id","unit":"","scale":"lineal","measured":false, ...} ] }

// scene/appearance.gs.json -- the trained appearance embedded in scene.glb
{ "role": "apariencia real entrenada con gsplat", "measured": false,
  "profile": "histora-gs-apariencia/1.0", "frame": "frame.ios_master", "n_primitives": 113827,
  "note": "f_dc_* = colour MEASURED PER TOOTH ... NOT measured=true: the optimiser moved,
           split and pruned the Gaussians, so there is no 1:1 correspondence with what was
           projected" }
\end{lstlisting}

\subsection{What a viewer must do, in order}\label{sec:vieweralgo}

\begin{enumerate}
  \item Read \code{manifest.json}; check \code{uos\_version}.
  \item Sort \code{assets} by \code{load\_priority} and fetch in that order. Several
  assets share \code{kind: mesh\_gs\_scene}; the one to fetch first is the glTF binary,
  which the manifest gives the lowest priority and which alone is enough to show the case.
  \item Load \code{scene/scene.glb} with any glTF 2.0 loader. Node 0 is already in the
  canonical frame; children are already positioned.
  \item For each child node carrying \code{extras.uos\_gs\_uri}: fetch the descriptor,
  \emph{then} the payload, \emph{then} apply the node matrix. Skip the layer if the
  descriptor is missing --- do not guess the column semantics.
  \item Read \code{views.json}, apply the default or fragment-selected view
  (\S\ref{sec:views}).
  \item Read \code{registrations[]} and mark every provisional one in the UI.
  \item If \code{derived/} is present and the deployment is permitted to show inference,
  load it \emph{as a separate toggleable layer} labelled with its regulatory status.
\end{enumerate}

% =============================================================================
\section{Volume}\label{sec:volume}

\paragraph{A slice's identity is not its file's hash.}\label{par:sliceid}
\code{Part.sha256} covers the whole file, so \textbf{any de-identification changes it} ---
and de-identification is the step every clinical flow takes, and the one \S\ref{sec:phi}
now requires. The traceability \S\ref{sec:external} promises --- ``whoever holds the series
can prove it is this case's, slice by slice'' --- broke at exactly the most common step.

DICOM already settled what identifies an instance: the \textbf{SOP Instance UID}
\code{(0008,0018)}, which survives de-identification when UIDs are retained, and the pixel
content \code{(7FE0,0010)}, which de-identification does not touch unless it is explicitly
cleaned. So \code{Part} additionally carries \code{sop\_instance\_uid} and
\code{pixel\_data\_sha256} (SHA-256 over the PixelData element's value, without the
header), and the \code{volume} asset carries \code{series\_instance\_uid} and
\code{study\_instance\_uid}.

The \code{parts[]} digest is computed over the UID and the pixel hash when every part
declares them, and falls back to name and file hash when they do not --- so an older
container still verifies. It therefore identifies the \textbf{clinical series} rather than
one serialisation of it. File names stay in \code{parts[]} as ordering data, which is what
\S\ref{sec:volume}'s argument for keeping them was about, but they leave identity.

Verification becomes two levels, reported separately: \textbf{identity} (UID + pixels) and
\textbf{exact bytes} (file hash). A de-identified file that preserves identity passes the
first and fails the second, and \emph{that is information, not an error}.

\paragraph{\code{value\_range} is measured, not skipped.}
It \MUST{} be computed \textbf{during the \code{parts[]} hash}, whenever the series is
available at ingestion --- which is always for the writer, by \S\ref{sec:external}: the
writer holds it, the recipient is the one who does not. The earlier justification for
leaving it null was that sweeping every pixel on every export is too expensive, and it does
not hold: the writer already reads every byte of every slice, for the file hash and for the
\code{pixel\_data\_sha256} above, so the minimum and maximum come out of a read that was
already paid for. Computing them anywhere else adds a whole extra pass over the volume for
a value that is free.

\code{null} is reserved for the one honest case: a container from an issuer that never had
the pixels.

\paragraph{A CBCT does not measure Hounsfield units.}
Grey values from a CBCT depend on the device, the field of view and the position inside the
volume; they are not comparable between scanners and not convertible to HU without a
calibration phantom. The volume sidecar therefore declares \code{calibrated\_hu}, true only
for a conventional CT. \textbf{Applying \code{rescale} to a CBCT does not yield Hounsfield
units}, and a reader that assumes it does will window on a scale that does not exist.

A CBCT series is a \textbf{directory asset} (\S\ref{sec:dirasset}): \code{kind: volume},
one \code{parts[]} entry per slice with its own hash, and, being an acquired original, a content address for its \code{uri} (\S\ref{sec:external}). The slices do not travel;
\code{parts[]} and the sidecar do.

\paragraph{What \code{parts[]} guarantees, and what a validator can check.}
Per-slice hashes let \textbf{whoever holds the series} prove it is the one this case was
built from, and prove it slice by slice, which is the difference between ``this series
does not match'' and ``slice 214 is corrupt''. The container itself cannot check a series
it does not hold, and \MUSTNOT{} claim to. Where a series \emph{is} carried (a container from another emitter, since ours never does) a validator \MUST{} check it in both directions: a slice present in the ZIP that
the manifest does not declare is as fatal as a declared slice that is missing. Either means
the series coming out is not the one claimed to have gone in. The reference test suite
builds such a container itself and removes a slice, adds a spurious one and alters one,
expecting three distinct errors.

\newpage
Slice names travel verbatim in \code{parts[]}. Renaming them to \code{IM0001.dcm} would
change the digest and describe a series nobody has, and the ordering of a series is clinical
data. Names produced by CBCT exports are equipment-generated (\code{3DSlice100.dcm}); a
vendor emitting patient identifiers in filenames is an ingestion problem, not a container
problem, by then the DICOM carries them in its tags anyway.

\subsection{The volume sidecar}

\code{sidecar\_uri} points at \code{volume/<id>.volume.json}. It exists so that a web
viewer does not need a complete DICOM parser (hundreds of files, transfer syntaxes,
per-slice rescale) merely to learn what the volume is. It matters more, not less, now
that the series is referenced: the sidecar is all a reader has until it resolves the series
elsewhere. It \MUST{} be read from the ingested bytes themselves, never from a description
computed somewhere else, and it \MUST{} travel even though the series does not.

\begin{longtable}{L{4.6cm} L{2.6cm} L{7.8cm}}
\toprule
\textbf{Field} & \textbf{Type} & \textbf{Description}\\
\midrule
\endhead
\code{frame} & string & The frame this volume is in. \textbf{Nothing else about alignment
appears here} --- the transform lives in \code{registrations[]}, and a system with two
homes for one transform ends up with two different transforms.\\
\code{dimensions} & int[3] & Columns, rows, slice count.\\
\code{spacing\_mm} & number[3] & $(x, y, z)$; $z$ from the measured inter-slice distance
when available, which overrides the declared \code{SliceThickness}.\\
\code{orientation} & number[3][3] & Rows from \code{ImageOrientationPatient}, third row
their cross product. \textbf{Read, never assumed} --- a volume with an assumed orientation
renders just as well, mirrored.\\
\code{origin\_mm} & number[3]\,|\,null & \code{ImagePositionPatient} of the first slice;
\code{null} with a warning if absent.\\
\code{rescale} & object & \code{\{slope, intercept\}}.\\
\code{value\_range} & number[2]\,|\,null & Rescaled $[\min,\max]$, from
\code{Smallest/LargestImagePixelValue} when the series declares them. \textbf{\code{null}
otherwise}, with a warning: sweeping every pixel of the series on each export is too
expensive, and inventing a plausible CBCT range would feed a fabricated number straight
into a viewer's windowing.\\
\code{pixel\_encoding} & string & e.g. \code{int16-le}. Byte order comes from the transfer
syntax, not from a convention --- ``almost all modern DICOM is little-endian'' is not
``all'', and a viewer that assumes reads good-looking garbage.\\
\code{modality} & string & DICOM \code{Modality}.\\
\code{transfer\_function\_presets} & string[] & Named presets a viewer may offer:
\code{cbct\_bone}, \code{cbct\_soft}, \code{cbct\_metal\_suppress}. These are names, not
data.\\
\bottomrule
\end{longtable}

% =============================================================================
\section{2D imaging and documents}\label{sec:image2d}

\code{kind: image2d} assets are clinical photographs and 2D radiographs. They are acquired
originals, so they are referenced by content address and not carried (\S\ref{sec:external});
what the manifest keeps is their identity, their \code{media\_type} and the optional
\code{projection} object, which says what kind of image it is and which teeth it targets:

\begin{lstlisting}[style=js]
"projection": { "type": "intraoral_frontal", "fdi_targets": ["11","12","13"] }
\end{lstlisting}

\paragraph{Normative: \code{fdi\_targets} empty means \emph{not recorded}.}
It does not mean \emph{none}. A loose photograph from a clinic folder does not say which
tooth it points at, and inferring it from the pixels is exactly the work the format
declines to assume has been done. A consumer \MUSTNOT{} read the empty array as a negative
assertion.

\vspace{6px}
Camera \code{pose} is explicitly optional in v0.2 and the reference implementation does not
emit it: computing it requires photo$\leftrightarrow$mesh fusion, which is measurable and
does not converge cheaply without calibration.

\code{kind: document} covers PDFs, reports and any other file the case declares whose type
has no better \code{kind} --- including the referenced scanner STL, which is a mesh but not
part of the renderable scene.

% =============================================================================
\section{The clinical layer}\label{sec:clinical}

\code{clinical/observations.json}, declared by extension \code{ash\_clinical}.

\paragraph{Why it exists.} v0.2 gives per-tooth clinical attributes no place: the FHIR
mapping (\S\ref{sec:fhir}) sends them to an \code{Observation} on an external server. The practical consequence is
that a standalone \code{.uos} cannot answer ``what does the report say about tooth 24'',
which is the question a clinician asks while looking at the model. So it is declared as an
extension, carried as \code{kind: document}, mapped to \code{Observation} in
\code{fhir\_map}, and written down as \emph{ours} so nobody mistakes it for ratified.

\paragraph{Why layer 1 and not \code{derived/}.} What travels is the \textbf{transcription
of a report a person signed}. That is clinical record. Putting it under \code{derived/}
would make it removable, and removing \code{derived/} would then strip the case of what the
report says, which is not what that operation means.

\paragraph{Normative: extraction provenance per value.} Transcription is layer 1, but the \emph{extraction} may well be inference, and every value declares which. \code{derivation} is \code{deterministic} (a pattern found it written),
\code{inferred} (a model proposed it), or \code{null}. \textbf{\code{null} means not
declared, not deterministic}, and a consumer \MUSTNOT{} treat an undeclared value as
reproducible. Silence never gets promoted to an assertion.

\newpage
\begin{lstlisting}[style=js]
{ "schema": "histora-clinical/2.0",
  "vocabulary": "ISO-3950 (FDI) for teeth; the finding vocabulary is closed",
  "regulatory": {
    "layer": 1, "default": true,
    "note": "THE FILE DEFAULT, not a claim about all of its content: every value in
             teeth[] declares its own layer and overrides this one",
    "confidence": "a value's confidence is NOT how sure the extractor was: it is the
                   weakest link in the chain that attached this value to THIS tooth, and
                   it is dominated by the segmenter's confidence in the FDI code" },
  "teeth": [
    { "fdi": "11",
      "findings": { "value": ["aparato_ortodoncico"],
                    "regulatory": {"layer": 1}, "derivation": "deterministic",
                    "confidence": 0.745, "agent": "report-agent@0.1.0",
                    "observed": "2026-08-17T15:39:38.921947+00:00" },
      "n_roots": { "value": 1, "regulatory": {"layer": 1}, ... },
      "color":   { "regulatory": {"layer": 2}, "derivation": "deterministic",
                   "derived_from": ["sha256:baeb31fa..."],
                   "value": { "space": "CIELAB",
                              "cervical": [74.89, 9.83, 24.06],
                              "middle":   [74.83, 6.70, 20.93],
                              "incisal":  [71.56, 6.77, 19.14],
                              "from_photo": "sha256:baeb31fa...", "n_pixels": 33508,
                              "measured": true,
                              "illumination_slope": [0.649, 0.577, 0.677],
                              "note": "measured per tooth; NOT a certified shade-guide
                                       value; flash falloff removed using the patient's
                                       own gingiva as reference, so teeth ARE comparable
                                       to each other; the absolute level is not, because
                                       the photograph carries no grey reference" } } } ] }
\end{lstlisting}

\paragraph{Regulatory layer per value, not per file.}
This file used to declare \code{layer: 1} once, for everything in it, on the grounds that
it is ``the transcription of a report a person signed''. Most of it is. The \code{color}
block is not: nobody signed a CIELAB triple per crown third with the flash falloff
regressed out --- the pipeline computes it from the photographs. The format had extraction
provenance per value (\code{derivation}, answering \emph{how} a value was obtained) and no
regulatory layer per value, so nothing answered \emph{who answers for it}. A reader who
deleted \code{derived/} kept computed measurements believing they kept a signed report.

Two rules, both of them \S\ref{sec:derived}'s carried inside the file:

\begin{itemize}
  \item A value with \code{derivation: "inferred"} \MUSTNOT{} declare \code{layer: 1}.
  What a model proposed is neither acquisition nor transcription, however right it turns
  out to be.
  \item A value declaring \code{layer: 3} inside \code{clinical/} is an \textbf{error}.
  Layer 3 lives only under \code{derived/}; that is the whole of what makes deleting the
  directory mean something.
\end{itemize}

The file-level \code{regulatory} stays, as the \textbf{default} for values that do not
declare their own, and says so in the object itself.

\paragraph{What \code{confidence} measures, said in the file.}
A value can carry \code{derivation: "deterministic"} and \code{confidence: 0.745} without
contradicting itself, and the file has to say why or every reader will read it as a
contradiction. The two answer different questions: \code{derivation} is how the value was
extracted --- a pattern found it written --- and \code{confidence} is how much of the
chain that attached it to \emph{this tooth} holds up. It is a minimum over that chain, and
in practice the segmenter's confidence in the FDI code dominates it. Deleting it where the
extraction was deterministic, which is the obvious-looking fix, would throw away the one
thing worth knowing: the report says this, and we are not certain which tooth it is about.

For the same reason \code{color} carries \textbf{no} \code{confidence}. That number is
about the report-to-tooth chain and says nothing about a colour measurement, which is
qualified by its own support: \code{n\_pixels}, \code{illumination\_slope},
\code{measured}.

Note the shape of that \code{note}: it states what the measurement supports (relative
comparison between teeth) and what it does not (absolute shade), instead of shipping a
number that looks like a shade match.

% =============================================================================
\section{{derived/} --- the inference plane}\label{sec:derived}

\textbf{The three hard rules}
\begin{enumerate}
  \item Everything a model produced lives \textbf{only} under \code{derived/}.
  \item Every such asset declares \code{regulatory.layer: 3}, and \emph{nothing outside}
  \code{derived/} may declare layer 3. A validator \MUST{} report an error in both
  directions.
  \item Every such asset has a sidecar declaring which model produced it, at which version,
  with which weights hash, over which source assets.
\end{enumerate}
And the consequence that makes them worth enforcing: \textbf{a \code{.uos} with no
\code{derived/} is still valid and complete}. Deleting the directory and its manifest
entries is a supported operation, requiring no re-export.

\subsection{Tooth segmentation}

\begin{longtable}{L{5cm} L{10.4cm}}
\toprule
\code{derived/seg\_teeth.bin} & FDI codes as \code{int16} little-endian, one per vertex of
the source asset, in the order the scene declares them. \code{0} = unassigned.\\
\code{derived/seg\_teeth.meta.json} & The sidecar below.\\
\bottomrule
\end{longtable}

\begin{lstlisting}[style=js]
{ "model": {"name": "histora-seg-teeth", "version": null, "weights_sha256": "ea0060d5..."},
  "source_assets": ["asset.scene"],
  "regulatory": {"layer": 3, "clearances": [
     {"jurisdiction": "EU", "regime": "MDR", "status": "investigational",
      "reference": null}]},
  "generated": "2026-08-28T02:07:17.244507+00:00",
  "encoding": { "dtype": "int16-le", "count": 112067,
                "indexes": "one code per vertex of `asset.scene`, in the same order the
                            scene declares them. 0 = unassigned.",
                "vocabulary": "ISO-3950 (FDI)" } }
\end{lstlisting}

\paragraph{Why not DICOM SEG or NRRD.} Both are designed for segmentation over a grid. This
one is not on a grid: it is a label \textbf{per primitive}, which is a third legitimate
shape and the only one that loses nothing. Inventing a grid to fit a known format would
resample the data to suit the wrapper.

\paragraph{\code{weights\_sha256: null} is normative.} It is \code{null} when unknown and
\MUSTNOT{} be filled with the hash of something else. The field exists so the inference can
be reproduced; a hash that is not the weights' turns ``I don't know'' into ``I do'' and is
indistinguishable from outside.

% =============================================================================
\section{views.json}\label{sec:views}

A view is a stored presentation state: ``what the clinician was looking at''. Views are
layer 1, cheap, versionable, and they are the natural mechanism for deep links
(\code{caso.uos\#view=view.pieza\_24} --- a \textbf{proposed} fragment syntax, since the
media type that would define it is not registered and no viewer implements it yet) and for
longitudinal comparison: the same view applied to \code{v1} and \code{v2}, side by side.
The example container stores 18.

\begin{lstlisting}[style=js]
{ "views": [
  { "id": "view.oclusal", "label": "Oclusal", "visit": "v1",
    "camera": { "position": [-11.852,-179.860, 31.622],
                "target":   [ -3.013,   0.415,   3.980],
                "up":       [ -0.256,   0.159,   0.953], "fov": 35.0 },
    "layers": { "mesh": {"visible": true, "opacity": 1.0} },
    "mpr": null, "clip_planes": null } ] }
\end{lstlisting}

\begin{longtable}{L{2.9cm} L{2.9cm} C{1.2cm} L{7.8cm}}
\toprule
\textbf{Field} & \textbf{Type} & \textbf{Req.} & \textbf{Description}\\
\midrule
\endhead
\code{id} & string & \MUST & Unique within the file. Conventions: \code{view.oclusal},
\code{view.frontal}, \code{view.vestibular\_derecha}, \code{view.vestibular\_izquierda},
\code{view.pieza\_<FDI>}.\\
\code{label} & string & \MUST & Display name.\\
\code{visit} & string & \MUST & \textbf{Must} name a declared visit; a validator errors
otherwise.\\
\code{camera} & object & \MUST & \code{position}, \code{target}, \code{up} (3 numbers each,
mm, canonical frame) and \code{fov} in degrees, default 35.\\
\code{layers} & map & --- & Per layer \code{\{visible: bool, opacity: 0..1\}}. Known keys:
\code{mesh}, \code{gs}, \code{volume}.\\
\code{mpr} & object\,|\,null & --- & \code{\{enabled, plane, window\}}. \\
\code{clip\_planes} & array\,|\,null & --- & Clipping planes.\\
\bottomrule
\end{longtable}

\paragraph{Normative: \code{mpr} and \code{clip\_planes} are null when there is no volume.}
They are volume controls. In a \code{.uos} that carries no volume, emitting them --- even
empty --- implies there is a plane to cut and a window to adjust. When a volume \emph{is}
present, omitting them is the opposite error: a viewer expects them.

\vspace{6px}
Four conventions worth copying, each with its reason:

\begin{itemize}
  \item The appearance layer defaults to \code{opacity: 0.85}, below 1.0, because it is
  reconstructed rather than measured and the viewer must let the mesh show through.
  \item The volume layer defaults to \textbf{off}. Composing translucent Gaussians with a
  translucent volume has no exact solution in two independent passes, so enabling it by
  default would store the view in the one state the format itself calls ill-defined.
  \item \code{fov} of 35$^\circ$: narrower flattens the arch, wider distorts its edges. It
  is the range in which a molar does not read as a cylinder.
  \item \code{time\_cursor\_s} is defined by the format for signal timelines and is not
  emitted by a writer that produces no signals --- a temporal cursor over a case with no
  time series points at nothing.
\end{itemize}

% =============================================================================
\section{Provenance, versioning and PHI}\label{sec:provenance}

\subsection{A \code{.uos} is logically append-only}

Modifying a case is \textbf{not} editing its \code{.uos}. It is writing a new version of
the manifest that points at the hash of the previous one. Assets are never touched in
place. That turns ``this case has changed'' into something the recipient can verify without
trusting the sender and without holding the intermediate versions.

\begin{longtable}{L{5.7cm} L{9.7cm}}
\toprule
\code{provenance.prev\_manifest\_sha256} & \textbf{Where the authority lives.} The SHA-256
of the bytes of the previous version's \code{manifest.json}.\\
\code{provenance.chain} & Path to the materialised chain; \MUST{} be
\code{"provenance/chain.json"} when present.\\
\bottomrule
\end{longtable}

\begin{lstlisting}[style=js]
{ "case_id": "urn:uuid:6a91c924-de59-5abc-94ec-8432c9f446ed",
  "links": [
    { "version": 1, "manifest_sha256": "f25b5a65...", "prev_manifest_sha256": null,
      "created": "2026-08-27T21:20:52.280181Z",
      "generator": {"name": "agentic-smart-health", "version": "0.4.0"},
      "assets": 19, "note": "18 view(s), 1 registration(s)" },
    { "version": 2, "manifest_sha256": "5f2e4a35...",
      "prev_manifest_sha256": "f25b5a65...", ... } ] }
\end{lstlisting}

\paragraph{What the chain does and does not guarantee.}
\code{chain.json} is a \emph{materialisation}: it makes the history legible at a glance and
walkable without opening every version. It adds no cryptographic guarantee --- it is not
covered by any hash in the manifest, and whoever rewrites the file breaks nothing
cryptographic. What a validator does detect, and what actually matters, is the chain and
the manifests \textbf{telling different stories}.

\paragraph{Hashing the manifest.} The hash is over the \textbf{exact bytes written into the
ZIP entry}. A writer therefore \MUST{} serialise once and hash those bytes; re-serialising
to compute the hash leaves it pointing at a second copy that merely resembles the first.
The reference implementation emits canonical JSON (stable key order, \code{indent=1}, nulls
retained, UTF-8).

\paragraph{Signatures.} \code{provenance/signatures/} is reserved. \textbf{v0.2 does not
specify signing.} What is missing is not code but a decision --- which key signs (the
issuing clinic, the platform, or both) and where it lives. Signing with an invented key
would produce a \code{.uos} that \emph{looks} signed, which is worse than one that declares
it is not. A validator that encounters the directory \MUST{} warn that the signatures are
present and unverified rather than ignore them silently.

\textbf{When it is specified, UOS \MUSTNOT{} define a signature scheme of its own.} The
open question is one of key infrastructure, and healthcare has answered it. A signature is
over the exact bytes of \code{manifest.json} --- the ones \S\ref{sec:provenance} already
requires be hashed --- as a detached \textbf{JWS} (JSON Web Signature, RFC 7515) or
\textbf{COSE Sign1}. Both admit Ed25519, so the current reservation is not lost; both carry
a key identifier (\code{kid}) pointing at the issuing clinic or the platform; both have
libraries in every language and a worked answer for key rotation. Multiple signatures ---
clinic \emph{and} platform --- then cost nothing: one entry per signer. The validator
verifies what it can (format, and the trust chain if given one) and keeps reporting the
rest as present and unverified, exactly as today. This is \MUST{} for v1.0.

% =============================================================================
\section{Extensions}\label{sec:extensions}

\paragraph{One prefix, and who owns it.}
Every extension this specification defines is named \code{histora\_}. It is the
\textbf{single} reserved prefix: a writer \MUSTNOT{} coin a second one under this format's
name, and an implementer adding their own extension \MUST{} use a prefix of their own.
Earlier drafts used \code{ash\_}, which was the summer-grant project's prefix rather than
anything anybody held --- the specification told writers not to squat on it without saying
who owned it, which is not a rule that can be followed. It is renamed, not aliased: the
format is unpublished, no second implementer has written against it, and the cost of this
change is the lowest it will ever be.

UOS builds on glTF, which has had \code{extensionsUsed} / \code{extensionsRequired} since
glTF 1.0 and carries them unchanged into 2.0: a reader opens the file, sees which extensions it carries, and
knows whether it can read it fully, partly, or not at all. v0.2 of the draft did not
inherit that mechanism at container level, and the practical consequence showed up
immediately: our extensions lived in their own directories and \textbf{a foreign reader
ignored them without noticing that it was ignoring them}. That turns an open format into
one only its author reads completely --- exactly what UOS exists to avoid.

\paragraph{And that gap is what the rest of this section closes.} With the mechanism
below, a reader that does not implement \code{ash\_clinical} still opens the case, still
renders it, and can say precisely what it left unread. Being open was never the problem;
being silent about what a reader was missing was.

\begin{longtable}{L{2.9cm} L{2.6cm} C{1.2cm} L{8.1cm}}
\toprule
\textbf{Field} & \textbf{Type} & \textbf{Req.} & \textbf{Description}\\
\midrule
\endhead
\code{name} & string & \MUST & Registry key; also the map key in \code{extensions}.\\
\code{version} & string & \MUST & Extension version, independent of \code{uos\_version}.\\
\code{uri} & string\,|\,null & --- & The asset that carries it, if any. If non-null it
\MUST{} be a declared asset \code{uri}.\\
\code{schema\_id} & string\,|\,null & --- & Identifier of its schema, e.g.
\code{histora-clinical/1.0}.\\
\code{description} & string & --- & What it adds and why the base format needed it.\\
\bottomrule
\end{longtable}

\paragraph{Normative: \code{used} versus \code{required} is the whole mechanism.}
\begin{itemize}
  \item \code{extensions\_used} --- ``this file contains this; ignore it if you must''. A
  reader that does not know it \MUST{} warn and \MUST{} still open the case.
  \item \code{extensions\_required} --- ``\textbf{do not open this file without
  understanding this}''. A reader that does not implement it \MUST{} refuse.
  \item Everything in \code{extensions\_used} \MUST{} be declared in \code{extensions};
  everything in \code{extensions\_required} \MUST{} also appear in \code{extensions\_used}.
  Violations are errors.
  \item An extension that only \emph{adds} information belongs in \code{used} and
  \textbf{never} in \code{required} --- otherwise a conformant viewer refuses a case whose
  geometry it could have shown perfectly.
\end{itemize}
Every extension the reference implementation emits is \code{used}; \code{required} is
empty, deliberately.

\subsection{Registered extensions}

\begin{longtable}{L{3.0cm} L{2.0cm} L{10.0cm}}
\toprule
\textbf{Name} & \textbf{Version} & \textbf{What it adds}\\
\midrule
\endhead
\code{ash\_clinical} & 1.0 & Per-tooth clinical attributes (pH, roots, canals, findings,
measured colour) with per-value extraction provenance. \S\ref{sec:clinical}. Carried at
\code{clinical/observations.json}.\\
\code{ash\_gs\_measured} & 1.0 & The \code{.gs.json} descriptor per Gaussian layer:
whether it is \emph{measured} or reconstructed, and the schema of its columns.
\S\ref{sec:gsjson}. The base format treats 3DGS as appearance in a reconstruction frame;
here there are fields fitted to CBCT density, in the patient frame, with error in the field's own unit. \textbf{Not in HU}: a CBCT is not calibrated in Hounsfield units (\S\ref{sec:volume}).\\
\code{ash\_reversible} & 1.0 & Declares that an arch mesh with per-vertex colour, FDI code
and a \code{measured} column can be \emph{regenerated} from \code{asset.scene}: geometry
and FDI from the mesh primitives (\code{extras.uos\_fdi}), colour from the
\code{KHR\_gaussian\_splatting} primitive, taking Gaussians that cover each vertex within
$3\sigma$ and weighting by alpha. \textbf{The mesh does not travel: it is reconstructed.}\\
\bottomrule
\end{longtable}

\subsection{Defining your own}

\begin{enumerate}
  \item Pick a prefixed name (\code{<vendor>\_<thing>}); do not squat on \code{ash\_} or
  \code{KHR\_}.
  \item Put its payload at a path of your own, declared as an asset like any other, with
  its hash. Extension payloads are not exempt from verification.
  \item Add an \code{Extension} entry and list it in \code{extensions\_used}.
  \item Put it in \code{extensions\_required} \textbf{only} if a reader that ignores it
  would show the user something \emph{wrong} --- not merely something incomplete.
  \item If it applies to inference output, it still obeys \S\ref{sec:derived}: under
  \code{derived/}, layer 3, with a model sidecar.
\end{enumerate}

% =============================================================================
\section{Conformance and validation}\label{sec:conformance}

\subsection{Conformance levels}

Levels are derived from what a container \emph{actually holds}, not from what it claims. A
writer does not choose them.

\begin{longtable}{L{2.8cm} L{4.4cm} L{7.8cm}}
\toprule
\textbf{Level} & \textbf{Requires asset kinds} & \textbf{Meaning}\\
\midrule
\endhead
\code{UOS-Core} & \code{mesh\_gs\_scene} & Manifest + a renderable scene. The floor.\\
\code{UOS-Vol} & \code{+ volume} & Plus a DICOM volume.\\
\code{UOS-Sig} & \code{+ signal} & Plus time-series signals. \emph{Not implemented in
v0.2.}\\
\code{UOS-Full} & \code{mesh\_gs\_scene + volume + signal} & All of the above.\\
\bottomrule
\end{longtable}

A container reports every level whose required set is a subset of the kinds present.

\paragraph{Levels count declared kinds, not carried bytes.}
An acquired original is referenced rather than carried (\S\ref{sec:external}), and it is
still a declared asset of its \code{kind}. A container whose CBCT is referenced therefore
reports \code{UOS-Vol} exactly like one that somehow held the slices: the level says the
case \emph{has} a volume and that this container knows how it relates to everything else,
which is what an implementer needs in order to decide whether it can open it. What a level does \MUSTNOT{} be read as is a custody claim. Custody is stated separately,
per asset, by \code{external} --- and a reader that conflates the two will promise a file
it does not have.

\paragraph{\code{UOS-Distributable}: an orthogonal profile.}
The levels answer whether a reader can \emph{open} the container. Nothing answered whether
it is in a condition to \emph{leave} the organisation that issued it, which is the question
asked immediately before attaching a case to an email. That profile is derived the same way
the levels are --- checked, not chosen --- and a container is \code{UOS-Distributable} when
all of the following hold:

\begin{itemize}
  \item \code{phi\_state} is not \code{identified};
  \item \code{deidentification} is present and consistent with the payload
  (\S\ref{sec:phi}) --- in particular the face rule;
  \item \code{purpose\_of\_use} is declared and contained in \code{consent.scope}
  (\S\ref{sec:consent});
  \item no asset declares layer 3 outside \code{derived/} (\S\ref{sec:derived});
  \item the container validates against its own manifest.
\end{itemize}

They are one profile and not five flags because separately they decide nothing: a container
with its purpose declared and a reconstructible face inside is no more sendable than one
with neither.

\textbf{Not being distributable is not an error.} It is the normal state of a case while it
lives inside the clinic, and a validator \MUSTNOT{} report it as one. What it \MUST{} do is
say so, and say what is missing, before somebody attaches the file to an email.

\subsection{The validation algorithm}\label{sec:algo}

A conforming validator performs these checks in this order. \textbf{E} = error,
\textbf{W} = warning.

\begin{longtable}{C{0.8cm} C{0.9cm} L{13.4cm}}
\toprule
\textbf{\#} & \textbf{Cls} & \textbf{Check}\\
\midrule
\endhead
1 & E & First ZIP entry is \code{manifest.json}. If not, stop.\\
2 & E & No entry uses a compression method other than STORE.\\
3 & W & The manifest validates against the distributed JSON Schema. A failure here is
a warning, not an error: what it detects is that our distributed schema has fallen behind the
contract, which is our publication problem, not the file's.\\
4 & E & The manifest parses against the contract (types, required fields, no unknown
keys).\\
5 & E & For each non-external asset: the \code{uri} exists in the ZIP; its SHA-256 and byte
length match.\\
6 & W & Referenced originals: \textbf{one} line per container naming how many there are
and which, not one warning per asset. Report the count separately so a caller can act on it
without parsing prose.\\
7 & E & For each directory asset: \code{parts} is non-empty; every declared part exists and
matches; \textbf{no undeclared file exists under the prefix}; the parts digest
(\S\ref{sec:dirasset}) equals the asset \code{sha256}; the byte sum matches.\\
8 & E & Every \code{sidecar\_uri} exists in the container.\\
9 & E & Provenance: if \code{chain} is declared it equals \code{provenance/chain.json}, the
entry exists, parses, its \code{case\_id} matches, and its links are consistent with
\code{prev\_manifest\_sha256} and with the hash of this manifest.\\
10 & W & \code{prev\_manifest\_sha256} present but no \code{chain}: history exists but
cannot be walked.\\
11 & E & \code{provenance/chain.json} present in the ZIP but not declared: a chain nobody
references cannot be verified.\\
12 & W & \code{provenance/signatures/} present: $n$ signatures that this validator does not
check.\\
13 & E & Views: each \code{visit} names a declared visit; ids are unique.\\
14 & W & No \code{views.json}: no stored views.\\
15 & E & Frame connectivity (\S\ref{sec:frames}).\\
16 & W & Each provisional registration (\S\ref{sec:manifest}).\\
17 & E & Regulatory: every asset under \code{derived/} has layer 3, and every layer-3 asset
is under \code{derived/}. This verifies where layer 3 is \emph{declared}.\\
17b & E & Layer content inside the payloads (\S\ref{sec:picking}): for every asset whose
\code{regulatory.layer} is not 3, a \code{.glb} declares neither \code{extras.uos\_fdi} nor
the \code{\_REGION\_ID} attribute on any primitive, and a \code{.ply} declares no
\code{region\_id} property. This verifies where layer 3 \emph{is}, which is the other half
--- and without it check 17 passes on a container whose layer-1 payloads are full of model
output.\\
17c & E & Layer 2 declares \code{derived\_from} (\S\ref{sec:planes}); layer-3 assets with
an empty \code{clearances} raise a warning; a registration whose \code{operator} begins
with \code{auto:} declares \code{regulatory}.\\
17d & E & Clinical layers (\S\ref{sec:clinical}): inside \code{clinical/}, no value
declares layer 3, no \code{inferred} value declares layer 1, and every layer-2 value
declares \code{derived\_from}.\\
17e & E/W & PHI and purpose (\S\ref{sec:phi}, \S\ref{sec:consent}): a non-\code{identified}
\code{phi\_state} carries \code{deidentification} (E); each \code{volume}/\code{image2d}
absent from \code{applied\_to} warns (W); a measured volume-derived layer without
\code{CleanRecognizableVisualFeatures} is an error (E); \code{date\_shift\_days} outside
\code{identified} is an error (E); \code{purpose\_of\_use} outside
\code{consent.scope} is an error (E), and its absence a warning (W).\\
17f & E & \code{derived/} self-description (\S\ref{sec:derived}): every asset there has a
non-null \code{sidecar\_uri}; the sidecar declares \code{model.name}, a non-empty
\code{source\_assets} and \code{encoding}; declared labels are ISO 3950 codes (permanent or
deciduous) or 0; and \code{derived\_from} resolves to declared asset ids.\\
17g & E & Every ZIP entry is declared by the manifest (\S\ref{sec:container}). The other
direction was checked and this one was not --- an undeclared file travels with no hash
accrediting it and no regulatory layer, which is the shape a leak would take.\\
17h & --- & An external asset's \code{uri} is a content address and equals its
\code{sha256} (\S\ref{sec:external}, ``checked in both directions''). \textbf{This is
enforced by the schema, not by the algorithm}: it is a model-level rule, so a manifest with
that inconsistency does not parse at all and check 2 reports it. Listed here because the
rule is real and readers look for it in this table --- adding it to the algorithm as well
would be code that can never run, which suggests a coverage that comes from somewhere
else.\\
17i & E & \code{KHR\_gaussian\_splatting}: \code{SCALE} $\geq 0$, \code{OPACITY} in
$[0,1]$, unit \code{ROTATION}, and complete spherical-harmonic degrees --- if degree $l$ is
present, every lower degree is present in full.\\
17j & E & \code{derived/seg\_*} indexes what it claims: byte length matches
\code{encoding.count}; \code{count} matches the \emph{source asset's} vertex count, not
only its own sidecar; every code in the payload is an ISO 3950 code or 0, read from the
bytes rather than from the sidecar's own list of labels; and where
\code{source\_positions\_sha256} is declared it matches the source's \code{POSITION}
accessor.\\
17m & E & A node carrying \code{extras.uos\_gs\_uri} has a \code{matrix} equal to the
\textbf{transpose} of the registration taking its frame to the canonical one
(\S\ref{sec:graph}), within tolerance. glTF stores matrices column-major and the manifest
row-major; confusing them does not fail, it places the cloud rotated and mirrored.\\
17k & E/W & Views: \code{mpr}/\code{clip\_planes} present iff a \code{volume} asset is
(E); \code{up} not parallel to the view direction and \code{fov} in a sane range (W).\\
17l & E & Chain order: \code{links[i].version} equals $i+1$ and \code{created} does not go
backwards. Checking the hashes without the order lets through a history whose links match
one by one and tell a different story.\\
18 & E & Extensions: \code{used} $\subseteq$ \code{extensions}; \code{required} $\subseteq$
\code{used}; every non-null extension \code{uri} is a declared asset.\\
19 & W & Extensions declared but unused; and one warning per required extension, since a
reader that does not implement it must not open the container.\\
20 & E & \code{canonical\_frame.units == "mm"}.\\
21 & --- & Report the conformance levels satisfied, whether the container is
\code{UOS-Distributable} and, when it is not, what is missing; plus the chain length and
the view count.\\
\bottomrule
\end{longtable}

\paragraph{The order of check 4 is normative.}
``No unknown keys'' \MUST{} be applied \emph{after} the version branch \S\ref{sec:versioning}
requires. A container declaring a higher \textbf{minor} is parsed leniently, ignoring and
naming the fields the reader does not know, and checking it against a schema with
\code{additionalProperties: false} first contradicts the compatibility a minor is defined
to promise --- in the one case that promise exists for. A validator reading a higher minor
therefore \MUST{} skip check 3 and say that it skipped it, rather than report every new
field as unknown.

\paragraph{Why check 3 is only a warning.} Validating against the distributed schema is not
redundant with validating against the implementation's types. The implementation checks
that a manifest fits \emph{our} contract; the schema shipped alongside it is the only artifact a
third party can check \emph{theirs} against. Running both is what guarantees the schema we publish
is the one we actually accept --- and if they drift, this check says so before an outside
implementer discovers it.

\paragraph{Reporting.} A validator \SHOULD{} report the chain length as \code{version} and
0 when no chain is declared. Zero means ``there is no history to walk'', not ``this is the
first version''; they are different facts.

% =============================================================================
\section{Recipes}\label{sec:cookbook}

Python examples use the reference package. Every one of them has a language-agnostic
equivalent stated alongside, because none of this requires that package.

\subsection{Open a case and verify it}\label{rec:open}

\begin{lstlisting}[style=py]
from pathlib import Path
from uos import read_manifest, validate

m = read_manifest(Path("case.uos"))        # errors unless manifest.json is entry #0
print(m.case_id, m.uos_version, m.phi_state)

report = validate(Path("case.uos"))        # the algorithm of section 12.3
print(report.valid, report.levels)         # e.g. True [UOS-Core, UOS-Vol]
print(report.distributable)                # may it leave the issuing organisation?
for e in report.errors:   print("ERROR  ", e)
for w in report.warnings: print("WARNING", w)
\end{lstlisting}

Without the package: open the ZIP, assert entry 0 is \code{manifest.json}, parse it,
validate against the distributed JSON Schema, then walk the checks in \S\ref{sec:algo}. No
step needs anything but a ZIP reader, a JSON parser and SHA-256.

\subsection{Render it}\label{rec:render}

Follow \S\ref{sec:vieweralgo} exactly. The minimum viable viewer is: read the manifest,
extract \code{scene/scene.glb}, hand it to a stock glTF loader, apply
\code{views.json}[0]. That already gives correct geometry, correct appearance and the
clinician's framing, and it is roughly thirty lines. Everything after that --- volume, MPR,
provisional badges, derived toggles --- is additive.

\subsection{Add an asset (e.g. a clinical photograph)}\label{rec:add}

\begin{lstlisting}[style=py]
from uos.contenedor import asset_de, escribe_uos
from uos.manifiesto import Clase

nuevo = asset_de(                       # hashes and measures the real file
    Path("/data/photo_09.jpg"), "images/img_009.jpg",
    id_="asset.img_009", kind=Clase.IMAGE2D, visit="v1",
    frame="frame.ios_master", media_type="image/jpeg",
)
m.assets.append(nuevo)
m.fhir_map["asset.img_009"] = RecursoFHIR(
    resource_type="Media", note="clinical photograph; Media is the non-DICOM image resource")
\end{lstlisting}

Checklist for \emph{any} new asset, in order:
\begin{enumerate}
  \item Choose a \code{kind} and a path consistent with \S\ref{sec:paths}.
  \item Hash the exact bytes you will store; set \code{sha256} and \code{bytes}.
  \item Name a \code{frame} that is reachable from the canonical one --- adding a new frame
  means also adding the \code{registrations[]} edge that connects it.
  \item Name a declared \code{visit}.
  \item If it came from a model, it goes under \code{derived/} with layer 3 and a sidecar,
  no exceptions (\S\ref{sec:derived}).
  \item If it came from another asset, populate \code{derived\_from}.
  \item Map it in \code{fhir\_map} if a connector should know what to create.
  \item Then write a \textbf{new version}, per \S\ref{rec:version}.
\end{enumerate}

\subsection{Add a view}

Append to \code{views.json} with a unique \code{id}, a declared \code{visit}, camera
vectors in millimetres in the canonical frame, and \code{mpr}/\code{clip\_planes} null
unless the container carries a volume. Views are cheap and rewriting them wholesale is
fine.

\subsection{Add a derived layer}

\begin{enumerate}
  \item Write the payload at \code{derived/<name>.<ext>} and a sidecar at
  \code{derived/<name>.meta.json} carrying \code{model} (name, version,
  \code{weights\_sha256} --- \code{null} if unknown, never a substitute hash),
  \code{source\_assets}, \code{regulatory} (layer 3, \code{clearances}),
  \code{generated}, and an \code{encoding} block that says \textbf{what the bytes index}.
  \item Declare both as assets with \code{regulatory.layer: 3}.
  \item Verify the removability invariant: delete \code{derived/} and its manifest entries,
  re-run the validator, and confirm the container is still valid. If it is not, you have
  put inference in the wrong plane.
\end{enumerate}

\subsection{Write version $N{+}1$}\label{rec:version}

\begin{lstlisting}[style=py]
# 1. serialise ONCE; the hash must be of the bytes that go into the ZIP
texto = manifiesto.json_canonico()
sha   = sha256_texto(texto)

# 2. link it to the previous version
manifiesto.provenance.prev_manifest_sha256 = sha_anterior
manifiesto.provenance.chain = "provenance/chain.json"
cadena = cadena_anterior.continua(Eslabon(
    version=len(cadena_anterior.links) + 1, manifest_sha256=sha,
    prev_manifest_sha256=sha_anterior, created=datetime.now(UTC),
    generator=manifiesto.generator, assets=len(manifiesto.assets), note="added img_009"))

# 3. write atomically; pass the SAME serialisation you hashed
escribe_uos(Path("case-v4.uos"), manifiesto, ficheros,
            directorios=dirs, extras={"provenance/chain.json": cadena.json_canonico()},
            json_manifiesto=texto)
\end{lstlisting}

Rules a writer \MUST{} honour here: assets are never edited in place; the file is written
to a temporary path and renamed, so a half-written container never sits where someone could
mistake it for the good one; and every \code{uri} in the manifest must have a file and
every supplied file must be declared --- a dangling reference is an error, not a gap.

\subsection{Declare an acquired original}\label{rec:external}

Every acquired original is declared this way; none is packed into the container
(\S\ref{sec:external}). Set \code{external: true}, put the content address
\code{sha256:<hash>} in \code{uri}, and keep \code{sha256}, \code{bytes},
\code{media\_type} and --- for a series --- the whole \code{parts[]} array. Referencing
gives up \emph{custody}, not traceability: whoever holds the series can still prove that
no slice is missing.

\MUSTNOT: a local filesystem path in \code{uri}. That is what would leak exactly what the
pseudonym exists to prevent.

\subsection{Regenerate the improved mesh}\label{rec:reversible}

This is what \code{ash\_reversible} declares, and the practical answer to ``why keep the
container if I already have the STL''. What comes back is not the file that went in: it is
that geometry with a per-vertex colour measured from the patient's own photographs and a
per-vertex FDI code, neither of which the ingested STL ever had.

\begin{enumerate}
  \item Read \code{scene/scene.glb}. Take \code{POSITION} and the index buffers of the
  \emph{mesh} primitives; the \code{extras.uos\_fdi} of each primitive gives the code for
  the vertices its faces touch.
  \item From the \code{KHR\_gaussian\_splatting} primitive take \code{POSITION},
  \code{SH\_DEGREE\_0\_COEF\_0}, \code{OPACITY} and \code{SCALE}.
  \item For each mesh vertex, take the Gaussians within $3\sigma$ --- the same cut-off the
  extension defines (\S\ref{sec:khr}) --- and average their colour weighted by
  \emph{opacity $\times$ Gaussian falloff}. That is the colour \emph{seen} at that point,
  which is not the colour of the nearest centre.
  \item Record, per vertex, whether it received colour at all. Write it out as a
  \code{measured} column.
\end{enumerate}

\begin{lstlisting}[style=py]
from export_agents.malla_mejorada import (
    color_desde_gaussianas, descriptor, escribe_ply, escribe_3mf, escribe_stl_viscam)

rgb, medido = color_desde_gaussianas(pos, centros, f_dc, opacidad, escalas)
meta = descriptor(pos, caras, rgb, medido, fdi,
                  origen="asset.scene", sha256_malla=sha, piezas_con_color=piezas)
escribe_ply(ruta.with_suffix(".ply"), pos, caras, rgb, fdi, meta["comentarios"], medido)
\end{lstlisting}

\paragraph{A note on that snippet.} It reaches into \code{export\_agents}, which is the
\emph{project's} pipeline and not this format's reference implementation --- which is why
its names are not English like the \code{uos} package's. Nothing here requires that
package: the recipe above states what a reader has to do, and the code is one way of doing
it.

\paragraph{Normative: the colour \MUST{} come from the Gaussian layer.} Not from any
intermediate painted mesh the producing pipeline happened to hold. That distinction is the
whole point: a recipient with the container and nothing else --- no repository, a year later
--- has to be able to obtain this. If the colour came from a file that does not travel, the
Gaussian layer would be contributing nothing and the chain would not close.

\paragraph{The \code{measured} column is not decoration.} Without it, ``I have no colour
here'' and ``I have the wrong colour here'' look identical. On the reference case
108{,}922 of 112{,}067 vertices take colour from their own tooth; of the remainder, 97
inherit it from the nearest measured vertex and 3{,}041 are painted with a fallback
gradient that is \textbf{not the patient's colour}. One byte per vertex turns that from a
doubt into a question anyone can answer without opening the \code{.uos}.

\paragraph{Ambient occlusion is deliberately not applied.} \code{\_AO} is a visualisation
factor. It has no place in a file somebody may print or take measurements on: darkening a
crown because there is a fissure beside it would put a shadow of ours inside the patient's
data.

\paragraph{Three formats, because none alone suffices --- and each says what it loses.}
\code{.ply} carries per-vertex colour, the FDI code and the \code{measured} flag, and any
mesh viewer opens it. \code{.3mf} is what slicers accept with colour and units inside.
\code{.stl} with per-face colour is a non-standard convention almost nothing reads, emitted
because there are laboratory workflows that accept nothing else --- it is the most widely
opened and the one that asserts least, and the descriptor says so rather than letting the
loss pass for absence.

\paragraph{Whole arch, no root, no per-tooth files.} Composing a scanner crown with a CBCT
root joins two surfaces measured by different machines across a 0.666\,mm registration, and
the seam shows. The scanner STL has no root, so giving it one makes it a different object:
reversibility asks for the same piece back, improved. Per-tooth files additionally depend on
a segmentation that does not currently separate crowns reliably (\S\ref{sec:limits}).

\subsection{Prepare a container for distribution}

Stripping inference is one of three steps, and doing only it produces a container that is
regulatorily clean and still legally undeclared.

\begin{enumerate}
  \item \textbf{Remove the inference plane.} Delete every entry under \code{derived/},
  delete the corresponding \code{assets[]} entries, re-serialise, and write a new version
  with a note saying what was removed.
  \item \textbf{Declare the purpose.} Set \code{purpose\_of\_use} to the act this copy is
  for, and check it against \code{subject.consent.scope} (\S\ref{sec:consent}).
  \item \textbf{Check the PHI claim against the payload} (\S\ref{sec:phi}). In particular,
  a container that still carries a \code{measured} Gaussian layer derived from a volume
  cannot claim \code{pseudonymized} unless recognisable visual features were cleaned before
  that layer was generated.
\end{enumerate}

Re-validate. The result is a valid, complete \code{.uos} at the same conformance levels.

% =============================================================================
\section{Versioning and compatibility}\label{sec:versioning}

\code{uos\_version} is the first field a reader checks, and this section says what checking
it buys. The format is versioned so that it can keep being developed without stranding the
implementations already written against it.

\subsection{What each kind of version change promises}

\begin{longtable}{L{3.0cm} L{12.2cm}}
\toprule
\textbf{Change} & \textbf{Guarantee}\\
\midrule
\endhead
\textbf{Minor} (0.2 $\to$ 0.3) & A reader written for the earlier version \MUST{} still be
able to open the container. A minor \MAY{} add optional fields, add extensions, and add
values to an open vocabulary. It \MUSTNOT{} add a required field, change what an existing
field means, or remove one.\\
\textbf{Major} (0.x $\to$ 1.0) & No compatibility is promised. A reader \MUST{} check
\code{uos\_version} and refuse what it does not implement rather than guess.\\
\bottomrule
\end{longtable}

\subsection{How that survives a strict schema}

\paragraph{The rule that makes minor compatibility real}
Each version publishes \textbf{its own} schema, and every schema sets
\code{additionalProperties: false}. That strictness is deliberate and is \emph{about that
version}: within a version a reader must refuse rather than guess, which is the premise of
\S\ref{sec:manifest}. 
\vspace{6px}
It is not a claim about later versions. So:
\begin{itemize}
  \item Reading a container whose \code{uos\_version} is \textbf{the reader's own}, an
  unknown field is an error.
  \item Reading a container whose \code{uos\_version} is a \textbf{higher minor}, an
  unknown field \MUST{} be ignored, and the reader \MUST{} report which version it read as
  which. Validating it against the older schema is not a conformance failure of the
  container.
\end{itemize}
This is why check 3 of \S\ref{sec:algo} is a warning and not an error: the contract decides
validity, and a schema left behind is the publisher's problem, not the file's.

\paragraph{Reading a higher minor, in practice.} A reader \MUST{} branch on the declared
version \emph{before} validating anything: a container it will refuse is not one it should
have parsed. The reference implementation does this in \code{uos.version}, and the lenient
path asks the validator which fields it rejected rather than keeping a second, permissive
copy of the model --- a duplicated model drifts from the real one at the first field
somebody adds, which is the same failure that makes this document's schema generated rather
than hand-written.

\paragraph{Normative: reading a higher minor and re-emitting it are mutually exclusive.}
A reader that ignores fields it does not understand \MUSTNOT{} then write version $N{+}1$
of that case. It would drop those fields, and the provenance chain would still verify: the
new version would be a cryptographically legitimate successor of the old one, with content
silently missing. That is worse than refusing, because nothing downstream can detect it.
Preserving unknown fields and emitting them back is the other way out, and it is more work
than is needed while no later version exists. So the rule is the narrow one: read it, do
not succeed it.

\subsection{Deprecation}

A field being withdrawn \MUST{} spend one full version marked deprecated before it goes:

\begin{enumerate}
  \item In the version that deprecates it, the field is still emitted, still valid, and the
  validator \MUST{} warn that it is on its way out and say what replaces it.
  \item In the next version it is removed from the schema.
\end{enumerate}

One version of notice is what lets an implementer migrate on their own schedule instead of
discovering the change when their reader breaks. The warning is also the written record of
\emph{why} the field went, which a diff of the schema does not carry.

\subsection{What is versioned alongside the format}

\begin{itemize}
  \item The manifest schema: \code{uos-manifest-<version>.schema.json}, whose \code{\$id}
  is pinned to a release tag so the identifier always names the same document
  (\S\ref{sec:artifacts}). A published tag \MUSTNOT{} be moved; a correction is a new one.
  \item Extensions, independently. \code{ash\_clinical} at 1.0 does not become 1.1 because
  UOS did (\S\ref{sec:extensions}).
  \item This document, whose version is on its title page.
\end{itemize}

% =============================================================================
\section{The conformance fixture}\label{sec:fixture}

A specification whose only executable definition lives inside its reference implementation
forces a second implementer to check their reader against \emph{their own output}, which is
checking nothing. The fixture is what makes that unnecessary.

It ships one \textbf{valid} container and several broken ones, each carrying \emph{a single}
defect, with the class of result each should produce written beside it in
\code{esperado.json}. An implementer runs their validator over the directory and compares.
Accepting one of the broken ones, or rejecting the valid one, tells them precisely what
they are missing.

\begin{longtable}{L{5.4cm} L{2.2cm} L{7.4cm}}
\toprule
\textbf{Case} & \textbf{Expects} & \textbf{What it exercises}\\
\midrule
\endhead
\code{valido.uos} & valid & The reference container.\\
\code{entrada-no-declarada.uos} & error & A file the manifest does not name: no hash
accrediting it, no regulatory layer (\S\ref{sec:container}).\\
\code{hash-que-no-cuadra.uos} & error & A carried asset whose declared \code{sha256} is not
its bytes'.\\
\code{manifiesto-no-es-primero.uos} & error & \code{manifest.json} is not the first
physical entry (\S\ref{sec:container}).\\
\code{capa-3-fuera-de-derived.uos} & error & Layer 3 outside \code{derived/}
(\S\ref{sec:derived}).\\
\code{registro-provisional.uos} & valid, warns & An automatic registration with no
\code{verified\_by}: the container is valid and a viewer must present it as provisional
(\S\ref{sec:manifest}).\\
\code{minor-superior-con-campo-nuevo.uos} & valid, warns & A higher minor with an unknown
field: it is ignored and \emph{named}, not rejected (\S\ref{sec:versioning}).\\
\code{comprimido.uos} & error & DEFLATE instead of STORE, so range access does not work
(\S\ref{sec:container}).\\
\bottomrule
\end{longtable}

\paragraph{Two properties it has on purpose.}
It is built \textbf{by running the real writer}, not by assembling manifests by hand: a
hand-built fixture tests what its author believes the writer emits, and this one tests what
the writer emits. And every case that patches the manifest also drops the provenance chain,
because \S\ref{sec:provenance} hashes the manifest's bytes --- a case that means to
exercise one rule and fails two exercises neither, since the implementer cannot tell which
one their reader rejected.

\textbf{All of its data is synthetic}, which is the condition for it being publishable at
all (\S\ref{sec:phi}). A fixture built on a real case would be useful once and
distributable never.

The fixture is regenerated and checked against this table on every run of the reference
implementation's test suite. A fixture nobody verifies is worse than none: if a case
stopped producing the result it advertises, an implementer would conclude their reader is
wrong when the fixture is.

\section{How this document changes}\label{sec:process}

A living specification has to say where changes are opened and who decides, or ``the
process of \S16'' names a process that does not exist --- which is what the previous
revision did.

\begin{longtable}{L{3.4cm} L{12.0cm}}
\toprule
\textbf{Question} & \textbf{Answer}\\
\midrule
\endhead
Who maintains it & The editor named on the cover, who holds the copyright. Version 0.2 was
written during the ANFAIA summer grant on a clinical case contributed by Histora;
maintenance is not tied to that term.\\
Where changes are opened & As issues and pull requests against the public repository that
hosts the reference implementation and this source. A change to the format and the change
to the implementation that realises it belong in the same request: the schema is derived
from the contract, so they cannot diverge silently, and a proposal that does not say how it
would be validated is not yet a proposal.\\
What a change must carry & The rule, in \MUST/\SHOULD/\MAY{} terms; where a validator would
check it and at which severity; and what it does to a container written under the previous
version (\S\ref{sec:versioning}).\\
Who decides & The editor, in the open, on the record of the request. This is a proposal
with one implementation, not a standard with a committee, and \S\ref{sec:limits} says so
plainly rather than implying otherwise.\\
When it is published & A version is published by tagging the repository, which is what
makes the schema's \code{\$id} resolve to that version and not to whatever the branch holds
today.\\
\bottomrule
\end{longtable}

\section{What v0.2 deliberately does not do}\label{sec:limits}

\begin{itemize}
  \item \textbf{It validates nothing clinically.} Everything inferred ships as
  \code{investigational}, on purpose.
  \item \textbf{Signals (\code{UOS-Sig}) are unimplemented.} The \code{kind} exists; no
  writer emits it and \code{time\_cursor\_s} in views is therefore never populated. When it
  is implemented, UOS \MUSTNOT{} define a time-series format of its own: DICOM already
  defines waveform objects for signals attached to a study, and outside DICOM there are
  open, widely supported formats for biomedical time series. Either is preferable to
  inventing one, and saying so here is what stops the first implementation from doing it.
  \item \textbf{Signatures are unspecified.} \S\ref{sec:provenance} names JWS or COSE
  Sign1 as the destination, so that the first implementation does not invent a scheme; what
  is undecided is which key signs and where it lives.
  \item \textbf{Camera pose for 2D imaging is not emitted.} \S\ref{sec:image2d}.
  \item \textbf{\code{COLOR\_0} is not emitted} on splat primitives, so non-extension
  viewers see an uncoloured point cloud. \S\ref{sec:khr}.
  \item \textbf{The Gaussian layers are monolithic.} STORE makes a container
  \emph{addressable} --- an HTTP client can fetch one asset without downloading the case ---
  and that is true for the manifest and for a 17.9\,MB GLB. It is not the same as
  \emph{streamable}, which the title suggests: \code{scene/field.ply} is 78\,MB in one file,
  so the range request is all-or-nothing, and on a clinic's poor connection that is not
  progressive. There is no level-of-detail and no per-region loading. A future version
  should evaluate splitting large layers into parts --- by FDI region or by spatial block ---
  declared as a directory asset with \code{parts[]}, reusing \S\ref{sec:external}'s existing
  mechanism exactly; a reader could then request only the parts falling inside the active
  view's frustum.
  \item \textbf{No partitioned scene is emitted under \code{derived/}.}
  \S\ref{sec:picking} defines the conforming way to keep per-tooth selection for a generic
  glTF reader --- a \emph{second} GLB, declared layer 3, with its model sidecar --- and the
  reference writer does not emit it. Two reasons, and the first is the one that decides it.
  \textbf{The reader it would serve barely exists}: reaching
  \code{derived/scene\_partitioned.glb} means knowing that a \code{.uos} is a ZIP, that
  \code{derived/} exists, and that an alternative scene lives inside it --- and a reader
  that knows all three already knows enough to join \code{derived/seg\_teeth.bin} by vertex
  index, which is 0.22\,MB against the 17.8\,MB the duplicate would cost. \textbf{And a
  duplicate with no arbiter is worse than the loss}: two copies of the same geometry, each
  individually valid, with nothing in the container declaring which one governs when they
  drift. A writer that has a concrete reader asking for it should emit it; one that does
  not, should not.
  \item \textbf{Every figure here was measured on a real patient's case}, which the project
  holds and does not publish. What is distributable is the \emph{conformance fixture}
  (\S\ref{sec:fixture}), built on synthetic data for exactly that reason.
  \item \textbf{The media type is unregistered.} \code{application/vnd.histora.uos}
  awaits an IANA registration in the \code{vnd.} tree, and \code{histora.dev} --- the
  namespace it names --- is not a domain anybody holds. The schema's \code{\$id} therefore
  resolves through the reference repository, not through that namespace.


  \item \textbf{It is not a standard.} It is a container proposal with a schema and a
  validator; the schema's \code{\$id} does not yet resolve to a published document.
\end{itemize}

% =============================================================================
\appendix
\section{Agent cheat sheet}\label{app:cheat}

Enough to act on a \code{.uos} without reading the rest of this document.

\begin{longtable}{L{4.6cm} L{10.8cm}}
\toprule
\textbf{Question} & \textbf{Answer}\\
\midrule
\endhead
Is this file a \code{.uos}? & ZIP whose first entry is \code{manifest.json} containing
\code{uos\_version}.\\
Where is everything? & \code{manifest.json} $\to$ \code{assets[]}. Nothing is discovered by
walking the ZIP; the manifest is exhaustive.\\
Can I trust a payload? & Re-hash it. \code{sha256} + \code{bytes} on every asset;
\code{parts[]} for directories.\\
Where is the geometry? & \code{scene/scene.glb}, stock glTF 2.0. Node 0 = canonical frame.\\
Where is the colour? & The \code{KHR\_gaussian\_splatting} primitive in that GLB.
$\mathrm{RGB} = \mathrm{SH}_0 \times 0.282094791 + 0.5$.\\
Which tooth is this? & Mesh primitive \code{extras.uos\_fdi} (string, FDI/ISO-3950), or
\code{\_REGION\_ID} per Gaussian, or \code{derived/seg\_teeth.bin} per vertex.\\
Where is the CBCT? & Directory asset under \code{volume/}, \code{kind: volume}, with a
\code{.volume.json} sidecar so you need no DICOM parser to inspect it.\\
How do two frames relate? & \code{registrations[]}, row-major, mm. Transpose for glTF.\\
Is this alignment trustworthy? & \code{rms\_error\_mm} + \code{verified\_by}. \code{null}
means nobody checked.\\
What did a model produce? & Everything under \code{derived/}, and only that. Layer 3.\\
What does the report say about tooth 24? & \code{clinical/observations.json},
\code{teeth[].fdi == "24"}.\\
Has this case changed? & \code{provenance.prev\_manifest\_sha256} and
\code{provenance/chain.json}.\\
Can I add data? & Yes: new asset + hash + frame + visit, then a \textbf{new version}. Never
edit in place.\\
Can I remove the AI output? & Yes: delete \code{derived/} and its manifest entries. The
container stays valid.\\
What must I never do? & Invent a hash; put a local path in an external \code{uri}; put a
patient name anywhere; re-order scene vertices; treat an empty \code{fdi\_targets} as ``no
teeth''; treat \code{derivation: null} as deterministic.\\
\bottomrule
\end{longtable}

\section{Measured figures}\label{app:figures}

Every number in this document comes from a real clinical case --- a maxillary arch, an
intraoral scan, a CBCT, nine photographs and three reports --- not from an example built
to illustrate it. They are grouped by what each one is evidence \emph{for}.

\subsection*{What the example container holds}

These are the figures a reader will meet on opening \code{another\_patient.uos}, and they
describe a container emitted exactly as \S\ref{sec:external} requires: the acquired
originals are referenced, everything UOS composes is inside.

\begin{longtable}{L{6.2cm} L{9.2cm}}
\toprule
\textbf{Item} & \textbf{Value}\\
\midrule
\endhead
ZIP entries & 12\\
Stored views & 18 (four anatomical, the rest one per annotated tooth)\\
Declared extensions & 3, all in \code{extensions\_used}, none required\\
Provenance chain & 3 versions\\
Mesh & 112{,}067 vertices, 220{,}085 triangles, 15 primitives (14 teeth + remainder)\\
Appearance layer & 113{,}827 Gaussians, SH degree 1, \code{srgb\_rec709\_display}\\
CBCT density field & 1{,}341{,}990 Gaussians, subsampled $3\times3\times1$ from
12{,}072{,}785\\
\bottomrule
\end{longtable}

\newpage
\subsection*{What has been measured, and what it supports}

\begin{longtable}{L{4.6cm} L{4.0cm} L{6.8cm}}
\toprule
\textbf{Measurement} & \textbf{Value} & \textbf{The claim it backs}\\
\midrule
\endhead
CBCT$\to$scanner registration & \code{icp\_surface}, RMS 0.666\,mm, \code{verified\_by}
null & That a registration carries its own error and its own verification status
(\S\ref{sec:manifest}) --- this one is provisional, a viewer must say so, and \textbf{the
validator says so too}: it emits a warning naming the registration, so the obligation does
not rest on the viewer remembering. It also now declares \code{regulatory} (layer 2, an
automatic computation) and \code{fit\_for: ["visualization"]} --- 0.666\,mm RMS is fine to
look at and does not license planning.\\

Mesh reversibility & $4.59\times10^{-6}$\,mm against a 0.1\,mm budget & That reversibility
means regenerating the mesh from what travels, not returning the ingested file
(\S\ref{sec:external}). The residual is the \code{float32} of the STL itself --- the order
of that format's own rounding unit, not a loss the conversion introduced.\\

Per-slice verification at scale & 397 slices, 419 entries, 0 errors, 0 warnings & That
\S\ref{sec:dirasset} works on a full series: a missing, altered or extra slice is caught
individually. Measured on a container written before originals stopped travelling
(\S\ref{sec:external}); the machinery it exercised is the same one that now emits
\code{parts[]} for a referenced series, so that whoever holds the CBCT can prove no slice
is missing.\\

Closed arch export & watertight, 12{,}025 closing faces, 11{,}936\,mm$^3$, 3.2\,s & That
the composed scene yields something the scan alone does not: a printable solid, and a
per-tooth file with a scanner-measured crown and a CBCT-reconstructed root.\\
\bottomrule
\end{longtable}

\section{References}\label{app:refs}

Every document this specification depends on or cites normatively. Web resources were
consulted on \textbf{2026-08}; where a link points at a branch rather than a tag, its
content may have moved since.

\subsection*{Normative dependencies}
A reader cannot implement UOS without these.

\begin{itemize}
  \item Khronos Group, \emph{glTF 2.0 Specification}.
  \url{https://registry.khronos.org/glTF/specs/2.0/glTF-2.0.html}
  \item Khronos Group, \emph{\code{KHR\_gaussian\_splatting}} --- \textbf{Release
  Candidate}, not yet ratified; ratification was announced for Q2 2026. Cited here as it
  stood on 2026-08-28.
  \url{https://github.com/KhronosGroup/glTF/tree/main/extensions/2.0/Khronos/KHR_gaussian_splatting}
  \item ISO/IEC 21320-1:2015, \emph{Document Container File --- Part 1: Core} --- the ZIP
  profile \S\ref{sec:container} requires.
  \item \emph{JSON Schema}, draft 2020-12 --- the manifest schema's dialect
  (\S\ref{sec:artifacts}). \url{https://json-schema.org/draft/2020-12/schema}
  \item RFC 2119, \emph{Key words for use in RFCs to Indicate Requirement Levels}, Bradner,
  March 1997 --- the conformance language of \S\ref{sec:intro}.
\end{itemize}

\subsection*{Referenced vocabularies and formats}
Needed to interpret what a container carries, not to open one.

\begin{itemize}
  \item NEMA, \emph{DICOM Standard}: PS3.10 (media storage and file format) and PS3.15
  Annex E.1 (attribute confidentiality profiles) --- \S\ref{sec:volume},
  \S\ref{sec:phi}.
  \item HL7, \emph{FHIR Release 4} (4.0.1), 2019: \code{ImagingStudy},
  \code{DocumentReference}, \code{Media}, \code{Observation} --- \S\ref{sec:fhir}.
  \item ISO 3950:2016, \emph{Dentistry --- Designation system for teeth and areas of the
  oral cavity} --- the FDI two-digit notation used throughout.
  \item RFC 3339, \emph{Date and Time on the Internet: Timestamps}, Klyne and Newman,
  July 2002 --- the format of every timestamp in the manifest.
  \item RFC 6838, \emph{Media Type Specifications and Registration Procedures}, Freed,
  Klensin and Hansen, January 2013 --- the registration template \S\ref{sec:container}
  follows.
  \item ASWF Color Interoperability Forum --- the naming behind
  \code{srgb\_rec709\_display} and \code{lin\_rec709\_display} (\S\ref{sec:khr}).
\end{itemize}

\subsection*{Precedents and recommended payload formats}
Cited as prior art or as the compression this specification points a writer towards
(\S\ref{sec:container}); neither is required.

\begin{itemize}
  \item Apple, \emph{USDZ File Format Specification} --- the uncompressed-ZIP precedent,
  including the alignment rule UOS does \emph{not} adopt.
  \url{https://openusd.org/release/spec_usdz.html}
  \item Niantic Labs, \emph{SPZ} --- a compressed interchange format for Gaussian splats.
  \url{https://github.com/nianticlabs/spz}
  \item Khronos Group, \emph{\code{KHR\_draco\_mesh\_compression}}.
  \url{https://github.com/KhronosGroup/glTF/tree/main/extensions/2.0/Khronos/KHR_draco_mesh_compression}
  \item Kerbl, Kopanas, Leimkühler and Drettakis, \emph{3D Gaussian Splatting for
  Real-Time Radiance Field Rendering}, ACM Transactions on Graphics 42(4), SIGGRAPH 2023
  --- the INRIA \code{.ply} conventions this format converts away from
  (\S\ref{sec:khr}).
\end{itemize}

\end{document}
